% Chapter written by unit TR-5 (plans/contraste/manuscripts-plan.md in CAOS_MANAGE).
% The tables and the numbers of this chapter are generated by tables/make_credit-capital.py from data/derived/C10 and the figures by
% figures/make_credit-capital.py; the \cc... macros come from tables/credit-capital-numbers.tex.
\makeatletter
\providecommand*{\input@path}{}
\g@addto@macro\input@path{{tables/}}
\makeatother
% Generated by tables/make_credit-capital.py from data/derived/C10: do not edit.
% One macro per number the credit-capital section quotes.
\newcommand{\ccAlpha}{5}
\newcommand{\ccAlphaQ}{17.506}
\newcommand{\ccAlphaQLow}{9.736}
\newcommand{\ccAnchorCells}{144}
\newcommand{\ccAnchorEs}{0.209404}
\newcommand{\ccAnchorHalfOk}{141}
\newcommand{\ccAnchorK}{0.070348}
\newcommand{\ccAnchorOff}{3}
\newcommand{\ccAnchorOffList}{other retail at a loss given default of 0.85 and a default probability of 0.1 percent (engine 21.0855, printed 21.08); other retail at a loss given default of 0.45 and a default probability of 1.3 percent (engine 50.7945, printed 50.80); residential mortgages at a loss given default of 0.45 and a default probability of 1.5 percent (engine 73.4441, printed 73.45)}
\newcommand{\ccAnchorVar}{0.176329}
\newcommand{\ccAnchorWorst}{0.0059}
\newcommand{\ccBiasMlWorstLong}{-0.007}
\newcommand{\ccBiasMlWorstShort}{-0.016}
\newcommand{\ccBiasMoWorstLong}{-0.023}
\newcommand{\ccBiasMoWorstShort}{-0.048}
\newcommand{\ccBoot}{500}
\newcommand{\ccBtLightCard}{red}
\newcommand{\ccBtLightCi}{green}
\newcommand{\ccBtLightCre}{green}
\newcommand{\ccBtLightOth}{red}
\newcommand{\ccBtLightRes}{red}
\newcommand{\ccCovMlLong}{0.940}
\newcommand{\ccCovMlMaxLong}{0.953}
\newcommand{\ccCovMlMaxShort}{0.955}
\newcommand{\ccCovMlMinLong}{0.932}
\newcommand{\ccCovMlMinRhoLong}{0.08}
\newcommand{\ccCovMlMinRhoShort}{0.02}
\newcommand{\ccCovMlMinSeLong}{1.6}
\newcommand{\ccCovMlMinSeShort}{3.4}
\newcommand{\ccCovMlMinShort}{0.912}
\newcommand{\ccCovMlShort}{0.948}
\newcommand{\ccCovMoLong}{0.870}
\newcommand{\ccCovMoMaxLong}{0.905}
\newcommand{\ccCovMoMaxShort}{0.843}
\newcommand{\ccCovMoMinLong}{0.807}
\newcommand{\ccCovMoMinShort}{0.655}
\newcommand{\ccCovMoShort}{0.815}
\newcommand{\ccCrpDCard}{0.00084}
\newcommand{\ccCrpDCi}{0.00126}
\newcommand{\ccCrpDCre}{0.00069}
\newcommand{\ccCrpDOth}{0.00100}
\newcommand{\ccCrpDRes}{0.00086}
\newcommand{\ccCrpOverGoneCard}{0.988}
\newcommand{\ccCrpOverGoneCi}{0.969}
\newcommand{\ccCrpOverGoneCre}{0.928}
\newcommand{\ccCrpOverGoneMax}{1.000}
\newcommand{\ccCrpOverGoneMin}{0.928}
\newcommand{\ccCrpOverGoneOth}{0.976}
\newcommand{\ccCrpOverGoneRes}{1.000}
\newcommand{\ccCrpPCard}{0.940}
\newcommand{\ccCrpPCi}{0.554}
\newcommand{\ccCrpPCre}{0.992}
\newcommand{\ccCrpPOth}{0.819}
\newcommand{\ccCrpPRes}{0.930}
\newcommand{\ccDelta}{4.8336}
\newcommand{\ccDeltaLow}{2.857}
\newcommand{\ccDeltaPaper}{4.83}
\newcommand{\ccEcLhpCard}{1.83}
\newcommand{\ccEcLhpCi}{2.40}
\newcommand{\ccEcLhpCre}{5.87}
\newcommand{\ccEcLhpOth}{0.53}
\newcommand{\ccEcLhpRes}{0.11}
\newcommand{\ccEffNamesCard}{1,667}
\newcommand{\ccEffNamesCi}{104}
\newcommand{\ccEffNamesCre}{50}
\newcommand{\ccEffNamesOth}{1,250}
\newcommand{\ccEffNamesRes}{595}
\newcommand{\ccElCard}{1.90}
\newcommand{\ccElCi}{0.41}
\newcommand{\ccElCre}{0.20}
\newcommand{\ccElOth}{0.44}
\newcommand{\ccElRes}{0.06}
\newcommand{\ccEngine}{1.1.0}
\newcommand{\ccEsCrpCard}{4.05}
\newcommand{\ccEsCrpCi}{3.51}
\newcommand{\ccEsCrpCre}{6.16}
\newcommand{\ccEsCrpOth}{1.18}
\newcommand{\ccEsCrpRes}{0.46}
\newcommand{\ccEsGmfCard}{3.81}
\newcommand{\ccEsGmfCi}{3.50}
\newcommand{\ccEsGmfCre}{6.21}
\newcommand{\ccEsGmfOth}{1.12}
\newcommand{\ccEsGmfRes}{0.45}
\newcommand{\ccEsGoneCard}{4.11}
\newcommand{\ccEsGoneCi}{3.68}
\newcommand{\ccEsGoneCre}{7.08}
\newcommand{\ccEsGoneOth}{1.21}
\newcommand{\ccEsGoneRes}{0.46}
\newcommand{\ccEsIsCard}{4.11}
\newcommand{\ccEsIsCi}{3.71}
\newcommand{\ccEsIsCre}{7.06}
\newcommand{\ccEsIsOth}{1.21}
\newcommand{\ccEsIsRes}{0.46}
\newcommand{\ccEsLhpCard}{3.96}
\newcommand{\ccEsLhpCi}{2.68}
\newcommand{\ccEsLhpCre}{6.01}
\newcommand{\ccEsLhpOth}{1.04}
\newcommand{\ccEsLhpRes}{0.20}
\newcommand{\ccEsStuCard}{17.13}
\newcommand{\ccEsStuCi}{16.00}
\newcommand{\ccEsStuCre}{18.97}
\newcommand{\ccEsStuOth}{13.23}
\newcommand{\ccEsStuRes}{8.23}
\newcommand{\ccEulerEsWorst}{\ensuremath{2.1\times10^{-14}}}
\newcommand{\ccEulerLargeAsyCard}{4.7}
\newcommand{\ccEulerLargeAsyCi}{25.5}
\newcommand{\ccEulerLargeAsyCre}{27.2}
\newcommand{\ccEulerLargeAsyOth}{9.9}
\newcommand{\ccEulerLargeAsyRes}{9.8}
\newcommand{\ccEulerLargeEsCard}{4.9}
\newcommand{\ccEulerLargeEsCi}{55.5}
\newcommand{\ccEulerLargeEsCre}{44.1}
\newcommand{\ccEulerLargeEsOth}{17.7}
\newcommand{\ccEulerLargeEsRes}{43.7}
\newcommand{\ccEulerRows}{60}
\newcommand{\ccEulerShareGap}{2.7}
\newcommand{\ccEulerTolVar}{0.01}
\newcommand{\ccEulerVarMaxCard}{0.0018}
\newcommand{\ccEulerVarMaxCi}{0.0015}
\newcommand{\ccEulerVarMaxCre}{0.0037}
\newcommand{\ccEulerVarMaxOth}{0.0013}
\newcommand{\ccEulerVarMaxRes}{0.0016}
\newcommand{\ccEulerVarWorst}{0.0037}
\newcommand{\ccEvFirst}{2007}
\newcommand{\ccEvLast}{2012}
\newcommand{\ccEvYears}{6}
\newcommand{\ccExcNFCard}{2}
\newcommand{\ccExcNFCi}{1}
\newcommand{\ccExcNFCre}{2}
\newcommand{\ccExcNFOth}{4}
\newcommand{\ccExcNFRes}{6}
\newcommand{\ccExcNNCard}{2}
\newcommand{\ccExcNNCi}{0}
\newcommand{\ccExcNNCre}{0}
\newcommand{\ccExcNNNCard}{2}
\newcommand{\ccExcNNNCi}{0}
\newcommand{\ccExcNNNCre}{0}
\newcommand{\ccExcNNNOth}{2}
\newcommand{\ccExcNNNRes}{5}
\newcommand{\ccExcNNOth}{3}
\newcommand{\ccExcNNRes}{5}
\newcommand{\ccExpNF}{0.3}
\newcommand{\ccExpNN}{0.06}
\newcommand{\ccExpNNN}{0.006}
\newcommand{\ccFitFirst}{1985}
\newcommand{\ccFitLast}{2006}
\newcommand{\ccFitMaxRateCard}{6.40}
\newcommand{\ccFitMaxRateCi}{1.80}
\newcommand{\ccFitMaxRateCre}{2.00}
\newcommand{\ccFitMaxRateOth}{1.44}
\newcommand{\ccFitMaxRateRes}{0.24}
\newcommand{\ccFitYears}{22}
\newcommand{\ccGaOverGoneCard}{0.990}
\newcommand{\ccGaOverGoneCi}{1.033}
\newcommand{\ccGaOverGoneCre}{1.192}
\newcommand{\ccGaOverGoneOth}{0.938}
\newcommand{\ccGaOverGoneRes}{0.824}
\newcommand{\ccGaRatioMax}{2.63}
\newcommand{\ccGaRatioMin}{1.15}
\newcommand{\ccGaShareCard}{2.5}
\newcommand{\ccGaShareCi}{42.8}
\newcommand{\ccGaShareCre}{43.7}
\newcommand{\ccGaShareMax}{87.8}
\newcommand{\ccGaShareMin}{2.5}
\newcommand{\ccGaShareOth}{9.0}
\newcommand{\ccGaShareRes}{87.8}
\newcommand{\ccGaTopCard}{0.095}
\newcommand{\ccGaTopCi}{0.983}
\newcommand{\ccGaTopCre}{1.982}
\newcommand{\ccGaTopOth}{0.087}
\newcommand{\ccGaTopRes}{0.156}
\newcommand{\ccGamma}{0.25}
\newcommand{\ccGmfOverGoneCard}{0.933}
\newcommand{\ccGmfOverGoneCi}{0.955}
\newcommand{\ccGmfOverGoneCre}{0.898}
\newcommand{\ccGmfOverGoneMax}{1.000}
\newcommand{\ccGmfOverGoneMin}{0.898}
\newcommand{\ccGmfOverGoneOth}{0.928}
\newcommand{\ccGmfOverGoneRes}{1.000}
\newcommand{\ccGoneOverLhpCard}{1.04}
\newcommand{\ccGoneOverLhpCi}{1.38}
\newcommand{\ccGoneOverLhpCre}{1.21}
\newcommand{\ccGoneOverLhpMax}{2.28}
\newcommand{\ccGoneOverLhpMin}{1.04}
\newcommand{\ccGoneOverLhpOth}{1.16}
\newcommand{\ccGoneOverLhpRes}{2.28}
\newcommand{\ccHhiCard}{0.0006}
\newcommand{\ccHhiCi}{0.0096}
\newcommand{\ccHhiCre}{0.0201}
\newcommand{\ccHhiOth}{0.0008}
\newcommand{\ccHhiRes}{0.0017}
\newcommand{\ccIsEsRatioCard}{237}
\newcommand{\ccIsEsRatioCi}{644}
\newcommand{\ccIsEsRatioCre}{1217}
\newcommand{\ccIsEsRatioMax}{1217}
\newcommand{\ccIsEsRatioMin}{237}
\newcommand{\ccIsEsRatioOth}{506}
\newcommand{\ccIsEsRatioRes}{995}
\newcommand{\ccIsEssCard}{0.005}
\newcommand{\ccIsEssCi}{0.017}
\newcommand{\ccIsEssCre}{0.077}
\newcommand{\ccIsEssOth}{0.006}
\newcommand{\ccIsEssRes}{1.708}
\newcommand{\ccIsOverGoneCard}{1.000}
\newcommand{\ccIsOverGoneCi}{1.008}
\newcommand{\ccIsOverGoneCre}{1.000}
\newcommand{\ccIsOverGoneMax}{1.022}
\newcommand{\ccIsOverGoneMin}{0.992}
\newcommand{\ccIsOverGoneOth}{0.992}
\newcommand{\ccIsOverGoneRes}{1.022}
\newcommand{\ccIsShiftCard}{-2.90}
\newcommand{\ccIsShiftCi}{-2.32}
\newcommand{\ccIsShiftCre}{-2.82}
\newcommand{\ccIsShiftOth}{-2.52}
\newcommand{\ccIsShiftRes}{-1.22}
\newcommand{\ccIsVarRatioCard}{120}
\newcommand{\ccIsVarRatioCi}{389}
\newcommand{\ccIsVarRatioCre}{369}
\newcommand{\ccIsVarRatioMax}{389}
\newcommand{\ccIsVarRatioMin}{60}
\newcommand{\ccIsVarRatioOth}{60}
\newcommand{\ccIsVarRatioRes}{93}
\newcommand{\ccKIrbCard}{3.92}
\newcommand{\ccKIrbCi}{7.14}
\newcommand{\ccKIrbCre}{6.68}
\newcommand{\ccKIrbCreCorp}{5.23}
\newcommand{\ccKIrbOth}{3.62}
\newcommand{\ccKIrbRes}{1.10}
\newcommand{\ccKstarCard}{3.87}
\newcommand{\ccKstarCi}{6.90}
\newcommand{\ccKstarCre}{6.45}
\newcommand{\ccKstarOth}{3.51}
\newcommand{\ccKstarRes}{1.09}
\newcommand{\ccLargeEadCard}{0.1}
\newcommand{\ccLargeEadCi}{2.0}
\newcommand{\ccLargeEadCre}{4.0}
\newcommand{\ccLargeEadOth}{0.2}
\newcommand{\ccLargeEadRes}{0.4}
\newcommand{\ccLargeNCard}{100}
\newcommand{\ccLargeNCi}{20}
\newcommand{\ccLargeNCre}{10}
\newcommand{\ccLargeNOth}{100}
\newcommand{\ccLargeNRes}{50}
\newcommand{\ccLgd}{0.45}
\newcommand{\ccLgdGamma}{0.25}
\newcommand{\ccLossUnit}{1,000}
\newcommand{\ccMatAdjCi}{1.269}
\newcommand{\ccMatAdjCre}{1.351}
\newcommand{\ccMcLightCard}{green}
\newcommand{\ccMcLightCi}{green}
\newcommand{\ccMcLightCre}{green}
\newcommand{\ccMcLightOth}{red}
\newcommand{\ccMcLightRes}{red}
\newcommand{\ccMcPCard}{0.205}
\newcommand{\ccMcPCi}{0.587}
\newcommand{\ccMcPCre}{0.131}
\newcommand{\ccMcPOth}{\ensuremath{1.1\times10^{-14}}}
\newcommand{\ccMcPRes}{\ensuremath{6.3\times10^{-121}}}
\newcommand{\ccMcZCard}{1.27}
\newcommand{\ccMcZCi}{-0.54}
\newcommand{\ccMcZCre}{-1.51}
\newcommand{\ccMcZOth}{7.73}
\newcommand{\ccMcZRes}{23.38}
\newcommand{\ccNFailedLong}{0}
\newcommand{\ccNFailedShort}{0}
\newcommand{\ccNGrid}{12}
\newcommand{\ccNMaxCard}{8,506}
\newcommand{\ccNMaxCi}{12,728}
\newcommand{\ccNMaxCre}{15,983}
\newcommand{\ccNMaxOth}{4,000}
\newcommand{\ccNMaxRes}{4,000}
\newcommand{\ccNamesCard}{3,000}
\newcommand{\ccNamesCi}{1,120}
\newcommand{\ccNamesCre}{550}
\newcommand{\ccNamesOth}{2,600}
\newcommand{\ccNamesRes}{2,250}
\newcommand{\ccNu}{5}
\newcommand{\ccNullRho}{0.12}
\newcommand{\ccOmegaCard}{0.125}
\newcommand{\ccOmegaCi}{1.055}
\newcommand{\ccOmegaCre}{2.062}
\newcommand{\ccOmegaOth}{0.500}
\newcommand{\ccOmegaRes}{1.064}
\newcommand{\ccPNFCard}{0.033}
\newcommand{\ccPNFCi}{0.265}
\newcommand{\ccPNFCre}{0.033}
\newcommand{\ccPNFOth}{\ensuremath{8.6\times10^{-5}}}
\newcommand{\ccPNFRes}{\ensuremath{1.6\times10^{-8}}}
\newcommand{\ccPNNCard}{0.001}
\newcommand{\ccPNNCi}{1.000}
\newcommand{\ccPNNCre}{1.000}
\newcommand{\ccPNNNCard}{\ensuremath{1.5\times10^{-5}}}
\newcommand{\ccPNNNCi}{1.000}
\newcommand{\ccPNNNCre}{1.000}
\newcommand{\ccPNNNOth}{\ensuremath{1.5\times10^{-5}}}
\newcommand{\ccPNNNRes}{\ensuremath{6.0\times10^{-15}}}
\newcommand{\ccPNNOth}{\ensuremath{2.0\times10^{-5}}}
\newcommand{\ccPNNRes}{\ensuremath{6.0\times10^{-10}}}
\newcommand{\ccPdMlCard}{4.228}
\newcommand{\ccPdMlCi}{0.910}
\newcommand{\ccPdMlCre}{0.436}
\newcommand{\ccPdMlHiCard}{4.717}
\newcommand{\ccPdMlHiCi}{1.297}
\newcommand{\ccPdMlHiCre}{1.336}
\newcommand{\ccPdMlHiOth}{1.110}
\newcommand{\ccPdMlHiRes}{0.173}
\newcommand{\ccPdMlLoCard}{3.815}
\newcommand{\ccPdMlLoCi}{0.680}
\newcommand{\ccPdMlLoCre}{0.188}
\newcommand{\ccPdMlLoOth}{0.866}
\newcommand{\ccPdMlLoRes}{0.116}
\newcommand{\ccPdMlOth}{0.975}
\newcommand{\ccPdMlRes}{0.139}
\newcommand{\ccPeakRateCard}{9.44}
\newcommand{\ccPeakRateCi}{2.30}
\newcommand{\ccPeakRateCre}{2.35}
\newcommand{\ccPeakRateOth}{3.06}
\newcommand{\ccPeakRateRes}{2.33}
\newcommand{\ccPeakYearCard}{2009}
\newcommand{\ccPeakYearCi}{2009}
\newcommand{\ccPeakYearCre}{2010}
\newcommand{\ccPeakYearOth}{2009}
\newcommand{\ccPeakYearRes}{2009}
\newcommand{\ccPowMlLongEight}{0.550}
\newcommand{\ccPowMlLongFour}{1.000}
\newcommand{\ccPowMlLongTen}{0.147}
\newcommand{\ccPowMlLongTwoFour}{0.968}
\newcommand{\ccPowMlShortEight}{0.318}
\newcommand{\ccPowMlShortFour}{0.990}
\newcommand{\ccPowMlShortTen}{0.152}
\newcommand{\ccPowMlShortTwoFour}{0.725}
\newcommand{\ccPowMoLongEight}{0.615}
\newcommand{\ccPowMoLongFour}{0.993}
\newcommand{\ccPowMoLongTen}{0.297}
\newcommand{\ccPowMoLongTwoFour}{0.305}
\newcommand{\ccPowMoShortEight}{0.585}
\newcommand{\ccPowMoShortFour}{0.985}
\newcommand{\ccPowMoShortTen}{0.350}
\newcommand{\ccPowMoShortTwoFour}{0.085}
\newcommand{\ccQNFCard}{6.07}
\newcommand{\ccQNFCi}{2.21}
\newcommand{\ccQNFCre}{1.85}
\newcommand{\ccQNFOth}{1.48}
\newcommand{\ccQNFRes}{0.24}
\newcommand{\ccQNNCard}{7.06}
\newcommand{\ccQNNCi}{3.34}
\newcommand{\ccQNNCre}{4.42}
\newcommand{\ccQNNNCard}{8.30}
\newcommand{\ccQNNNCi}{5.11}
\newcommand{\ccQNNNCre}{10.09}
\newcommand{\ccQNNNOth}{2.15}
\newcommand{\ccQNNNRes}{0.40}
\newcommand{\ccQNNOth}{1.77}
\newcommand{\ccQNNRes}{0.30}
\newcommand{\ccRatioCard}{0.47}
\newcommand{\ccRatioCi}{0.34}
\newcommand{\ccRatioCre}{0.88}
\newcommand{\ccRatioCreCorp}{1.12}
\newcommand{\ccRatioMax}{0.88}
\newcommand{\ccRatioMin}{0.10}
\newcommand{\ccRatioOth}{0.15}
\newcommand{\ccRatioRes}{0.10}
\newcommand{\ccRegScale}{1.060}
\newcommand{\ccRhoGridHi}{0.24}
\newcommand{\ccRhoGridLo}{0.02}
\newcommand{\ccRhoIrbCard}{0.040}
\newcommand{\ccRhoIrbCi}{0.196}
\newcommand{\ccRhoIrbCorp}{0.196}
\newcommand{\ccRhoIrbCre}{0.265}
\newcommand{\ccRhoIrbCreCorp}{0.216}
\newcommand{\ccRhoIrbOth}{0.122}
\newcommand{\ccRhoIrbRes}{0.150}
\newcommand{\ccRhoMlCard}{0.013}
\newcommand{\ccRhoMlCi}{0.064}
\newcommand{\ccRhoMlCre}{0.238}
\newcommand{\ccRhoMlHiCard}{0.024}
\newcommand{\ccRhoMlHiCi}{0.116}
\newcommand{\ccRhoMlHiCre}{0.407}
\newcommand{\ccRhoMlHiOth}{0.021}
\newcommand{\ccRhoMlHiRes}{0.028}
\newcommand{\ccRhoMlLoCard}{0.007}
\newcommand{\ccRhoMlLoCi}{0.038}
\newcommand{\ccRhoMlLoCre}{0.144}
\newcommand{\ccRhoMlLoOth}{0.006}
\newcommand{\ccRhoMlLoRes}{0.007}
\newcommand{\ccRhoMlOth}{0.011}
\newcommand{\ccRhoMlRes}{0.013}
\newcommand{\ccRhoMoCard}{0.012}
\newcommand{\ccRhoMoCi}{0.043}
\newcommand{\ccRhoMoCre}{0.156}
\newcommand{\ccRhoMoHiCard}{0.022}
\newcommand{\ccRhoMoHiCi}{0.076}
\newcommand{\ccRhoMoHiCre}{0.229}
\newcommand{\ccRhoMoHiOth}{0.018}
\newcommand{\ccRhoMoHiRes}{0.023}
\newcommand{\ccRhoMoLoCard}{0.006}
\newcommand{\ccRhoMoLoCi}{0.019}
\newcommand{\ccRhoMoLoCre}{0.052}
\newcommand{\ccRhoMoLoOth}{0.005}
\newcommand{\ccRhoMoLoRes}{0.005}
\newcommand{\ccRhoMoOth}{0.010}
\newcommand{\ccRhoMoRes}{0.013}
\newcommand{\ccRhoWidthCard}{3.2}
\newcommand{\ccRhoWidthCi}{3.0}
\newcommand{\ccRhoWidthCre}{2.8}
\newcommand{\ccRhoWidthMax}{4.0}
\newcommand{\ccRhoWidthMin}{2.8}
\newcommand{\ccRhoWidthOth}{3.3}
\newcommand{\ccRhoWidthRes}{4.0}
\newcommand{\ccScen}{1,000,000}
\newcommand{\ccScenIs}{100,000}
\newcommand{\ccSectorCorr}{0.5}
\newcommand{\ccSeedIs}{4004}
\newcommand{\ccSeedMf}{4002}
\newcommand{\ccSeedOne}{4001}
\newcommand{\ccSeedStu}{4003}
\newcommand{\ccSerLong}{40}
\newcommand{\ccSerShort}{22}
\newcommand{\ccSigmaCard}{0.0590}
\newcommand{\ccSigmaCi}{0.5521}
\newcommand{\ccSigmaCre}{4.5370}
\newcommand{\ccSigmaOth}{0.0805}
\newcommand{\ccSigmaRes}{0.1464}
\newcommand{\ccSimPd}{2}
\newcommand{\ccSimRep}{400}
\newcommand{\ccSimSeed}{7001}
\newcommand{\ccSizeMlLong}{0.060}
\newcommand{\ccSizeMlLongHi}{0.088}
\newcommand{\ccSizeMlLongLo}{0.041}
\newcommand{\ccSizeMlShort}{0.052}
\newcommand{\ccSizeMlShortHi}{0.079}
\newcommand{\ccSizeMlShortLo}{0.035}
\newcommand{\ccSizeMoLong}{0.130}
\newcommand{\ccSizeMoLongHi}{0.167}
\newcommand{\ccSizeMoLongLo}{0.101}
\newcommand{\ccSizeMoRatioLong}{2.6}
\newcommand{\ccSizeMoRatioShort}{3.7}
\newcommand{\ccSizeMoShort}{0.185}
\newcommand{\ccSizeMoShortHi}{0.226}
\newcommand{\ccSizeMoShortLo}{0.150}
\newcommand{\ccStuOverGoneCard}{4.0}
\newcommand{\ccStuOverGoneCi}{4.1}
\newcommand{\ccStuOverGoneCre}{2.6}
\newcommand{\ccStuOverGoneMax}{14.8}
\newcommand{\ccStuOverGoneMin}{2.6}
\newcommand{\ccStuOverGoneOth}{10.0}
\newcommand{\ccStuOverGoneRes}{14.8}
\newcommand{\ccTotalEad}{100}
\newcommand{\ccVarCrpCard}{3.83}
\newcommand{\ccVarCrpCi}{3.08}
\newcommand{\ccVarCrpCre}{5.08}
\newcommand{\ccVarCrpOth}{1.10}
\newcommand{\ccVarCrpRes}{0.41}
\newcommand{\ccVarGaCard}{3.83}
\newcommand{\ccVarGaCi}{3.28}
\newcommand{\ccVarGaCre}{6.52}
\newcommand{\ccVarGaOth}{1.05}
\newcommand{\ccVarGaRes}{0.33}
\newcommand{\ccVarGmfCard}{3.61}
\newcommand{\ccVarGmfCi}{3.03}
\newcommand{\ccVarGmfCre}{4.91}
\newcommand{\ccVarGmfOth}{1.04}
\newcommand{\ccVarGmfRes}{0.41}
\newcommand{\ccVarGoneCard}{3.87}
\newcommand{\ccVarGoneCi}{3.18}
\newcommand{\ccVarGoneCre}{5.47}
\newcommand{\ccVarGoneOth}{1.12}
\newcommand{\ccVarGoneRes}{0.41}
\newcommand{\ccVarIsCard}{3.87}
\newcommand{\ccVarIsCi}{3.20}
\newcommand{\ccVarIsCre}{5.47}
\newcommand{\ccVarIsOth}{1.12}
\newcommand{\ccVarIsRes}{0.41}
\newcommand{\ccVarLhpCard}{3.74}
\newcommand{\ccVarLhpCi}{2.30}
\newcommand{\ccVarLhpCre}{4.54}
\newcommand{\ccVarLhpOth}{0.97}
\newcommand{\ccVarLhpRes}{0.18}
\newcommand{\ccVarStuCard}{15.60}
\newcommand{\ccVarStuCi}{13.10}
\newcommand{\ccVarStuCre}{14.38}
\newcommand{\ccVarStuOth}{11.20}
\newcommand{\ccVarStuRes}{5.98}
\newcommand{\ccXi}{0.25}
\newcommand{\ccYearsCard}{22}
\newcommand{\ccYearsCi}{22}
\newcommand{\ccYearsCre}{16}
\newcommand{\ccYearsOth}{22}
\newcommand{\ccYearsRes}{16}
\newcommand{\ccYrsNNCard}{2009, 2010}
\newcommand{\ccYrsNNCi}{none}
\newcommand{\ccYrsNNCre}{none}
\newcommand{\ccYrsNNOth}{2008, 2009, 2010}
\newcommand{\ccYrsNNRes}{2008, 2009, 2010, 2011, 2012}

\emergencystretch=3em
\providecommand{\EL}{\mathrm{EL}}
\providecommand{\ccfn}[1]{\nolinkurl{#1}}

\section{Methods: credit capital and the portfolio loss distribution}
\label{sec:credit-capital}

Credit risk is the loss a bank bears when borrowers fail to pay. For one loan the expected loss is the product of the probability of
default, the loss given default and the exposure at default. The capital a bank holds against credit risk does not cover that average:
it covers the distance between the expected loss and a high quantile of the loss of the whole portfolio in one year. That distance is
set less by the average quality of the loans than by how strongly their defaults move together. A book of independent borrowers has a
narrow loss distribution; a book whose borrowers share one economic cycle has a long upper tail. The quantity that carries the
co-movement, the asset correlation, is also the hardest to observe, because a bank sees one default rate per year while the tail that
decides the capital is a one-in-a-thousand-year event \citep{credit-capital:vasicek2002,credit-capital:gordy2003}.

The regulation that governs the capital of a bank with supervisory approval is the internal ratings-based (IRB) approach of the Basel
Framework. For each exposure class it fixes a risk-weight function that turns the bank's estimates of the default probability and of
the loss given default into a capital requirement per unit of exposure at the 99.9 percent confidence level, with a correlation fixed
by the class \citep{credit-capital:basel-cre31,credit-capital:basel-cre32}. The function is the one-factor model of Vasicek
evaluated for an infinitely large portfolio \citep{credit-capital:basel2005}, and its justification as a portfolio-invariant capital
rule is the asymptotic single risk factor model of Gordy \citep{credit-capital:gordy2003}. The chapters in force from 1 January 2023
follow the final Basel III reforms \citep{credit-capital:basel-d424}; the European Union applies them through Article 153 of Regulation
(EU) 2024/1623 from 1 January 2025 \citep{credit-capital:crr3}; and the Basel II text, which multiplied the IRB risk-weighted assets by
1.06, remains the reference of older figures \citep{credit-capital:basel-cre30}. In Chile the Comisi\'on para el Mercado Financiero sets
the rules for internal methodologies in chapter 21-6 of its Recopilaci\'on Actualizada de Normas, with correlations of its own
\citep{credit-capital:cmf-ran21-6}, and its normative report of August 2026 states that no Chilean bank uses them for capital
\citep{credit-capital:cmf-informe2026}.

A portfolio-invariant formula cannot see the size distribution of the exposures, the concentration of the book in sectors, or the
dependence between defaults far in the tail. Those effects are the object of the granularity adjustment of
\citet{credit-capital:gordylutkebohmert2013}, of the studies of concentration risk
\citep{credit-capital:bcbs-wp15,credit-capital:dullmannmasschelein2007} and of the work on the copula that links the defaults
\citep{credit-capital:li2000,credit-capital:freymcneil2003}. A validator asked whether the capital of a loan book is adequate therefore
has three questions: which correlation the formula assumes and whether the loss history supports it, which tail the model has to
survive, and how much the portfolio's own shape and the family of the model change the capital.

One case of the product, C10, asks those questions on public data. The yearly net charge-off rates of five loan types (commercial and
industrial, commercial real estate, residential real estate, credit cards and other consumer) come from the charge-off release of the
Federal Reserve Board, all commercial banks, seasonally adjusted, as the mean of the four quarterly rates of each year
\citep{credit-capital:fed-chargeoff}. A Vasicek distribution is fitted on \ccFitFirst{} to \ccFitLast{} and the realised rates of
\ccEvFirst{} to \ccEvLast{} are placed in it, so the tail test is out of sample. For each loan type an illustrative portfolio of
\ccTotalEad{} million USD is built from the fitted default probability and correlation and is run through six models of increasing
realism: the large-portfolio formula, the Gaussian copula with one factor and with two sector factors, the Student-t copula,
CreditRisk+, the granularity adjustment and importance sampling. A simulator with a known correlation measures how often the two
correlation estimators are right. The charge-off rate is a net loss rate and not a default frequency, the series are those of all banks
and not of a portfolio, and the portfolios are constructed; the case therefore draws no statement about any bank. The data are
\ccFitYears{} yearly points at most, which is the reason the correlation is reported with an interval wherever it appears.

The methods are developed in two subsections. The first (\cref{sec:m16}) states the Vasicek loss distribution of the large homogeneous
portfolio, the IRB capital requirement and the calibration of the correlation from yearly loss rates. The second (\cref{sec:m17})
states the models of a finite portfolio, the measures they share, the validation tests that judge them and the allocation of the
capital to the names. The results of all methods on all loan types follow in \cref{sec:results-credit-capital}, and
\cref{tab:cc-cases} lists the five loan types with their series, their IRB class and the shape of their portfolio.

\begin{table}[htbp]
\centering
\small
\caption{The five loan types of case C10: the column of the Federal Reserve charge-off table that supplies the yearly loss rates, the nearest class of the Basel framework with its paragraph, the number of years the fit window holds, and the three groups of the illustrative portfolio (large, mid-size and small names, each shown as number of names times exposure of each name) with the effective number of names $1/\mathrm{HHI}$.}
\label{tab:cc-cases}
\begin{tabular}{>{\raggedright\arraybackslash}p{2.4cm}>{\raggedright\arraybackslash}p{2.2cm}>{\raggedright\arraybackslash}p{2.3cm}>{\raggedleft\arraybackslash}p{0.8cm}>{\raggedright\arraybackslash}p{3.6cm}>{\raggedleft\arraybackslash}p{1.4cm}}
\toprule
Loan type & Board series & IRB class (paragraph) & Fit years & Names $\times$ exposure (USD) & Effective names \\
\midrule
Commercial and industrial & C\&I loans (column CI) & corporate (CRE31.5) & 22 & $20\times 2{,}000$k / $100\times 400$k / $1000\times 20$k & 104 \\
Commercial real estate & Commercial (column RECI) & hvcre (CRE31.11) & 16 & $10\times 4{,}000$k / $40\times 1{,}000$k / $500\times 40$k & 50 \\
Residential real estate & Residential (column RERES) & residential mortgage (CRE31.14) & 16 & $50\times 400$k / $200\times 200$k / $2000\times 20$k & 595 \\
Credit cards & Credit cards (column CONCC) & qrre (CRE31.15) & 22 & $100\times 100$k / $400\times 100$k / $2500\times 20$k & 1,667 \\
Other consumer & Other (column CONOTH) & other retail (CRE31.16) & 22 & $100\times 200$k / $500\times 80$k / $2000\times 20$k & 1,250 \\
\bottomrule
\end{tabular}
\end{table}


\subsection{The Vasicek loss distribution and the IRB capital requirement}
\label{sec:m16}

\subsubsection{The single-factor model and the large homogeneous portfolio}

Let $\Phi$ be the standard normal distribution function and $\Phi^{-1}$ its inverse. In the model of
\citet{credit-capital:vasicek2002} (the notes are collected in \citealp{credit-capital:vasicek1987,credit-capital:vasicek1991}) the
latent asset return of obligor $i$ is
\begin{equation}
A_i=\sqrt{\rho}\,Z+\sqrt{1-\rho}\,\varepsilon_i,\qquad Z,\varepsilon_i\ \text{independent standard normal},
\label{eq:cc-asset}
\end{equation}
where $Z$ is the systematic factor shared by all obligors, $\varepsilon_i$ is the idiosyncratic factor of obligor $i$ and $\rho\in(0,1)$
is the asset correlation. Obligor $i$ defaults when $A_i<\Phi^{-1}(\PD)$, so the marginal default probability is $\PD$. Given $Z=z$ the
defaults are independent with the conditional default probability
\begin{equation}
p(z)=\Phi\!\left(\frac{\Phi^{-1}(\PD)-\sqrt{\rho}\,z}{\sqrt{1-\rho}}\right),
\label{eq:cc-condpd}
\end{equation}
which decreases in $z$: a low value of the factor is a bad year. When the number of obligors grows without bound with equal exposure
shares, the default fraction of the portfolio converges to $p(Z)$ (the large homogeneous portfolio, LHP). Because $p$ is monotone, the
event $\{p(Z)\le x\}$ equals $\{Z\ge z_x\}$ with $z_x=\big(\Phi^{-1}(\PD)-\sqrt{1-\rho}\,\Phi^{-1}(x)\big)/\sqrt{\rho}$, and therefore
\begin{equation}
\begin{aligned}
\Pr(L_\infty\le x)&=\Phi\!\left(\frac{\sqrt{1-\rho}\,\Phi^{-1}(x)-\Phi^{-1}(\PD)}{\sqrt{\rho}}\right),\\
\VaR_q(L_\infty)&=\Phi\!\left(\frac{\Phi^{-1}(\PD)+\sqrt{\rho}\,\Phi^{-1}(q)}{\sqrt{1-\rho}}\right),
\end{aligned}
\label{eq:cc-lhp}
\end{equation}
the second expression being the inverse of the first at $x=\VaR_q$ \citep{credit-capital:vasicek2002}. Differentiating the first with
$\mathrm{d}\Phi^{-1}(x)/\mathrm{d}x=1/\varphi(\Phi^{-1}(x))$, with $\varphi$ the standard normal density, gives the density of the default
fraction
\begin{equation}
f(x;\PD,\rho)=\sqrt{\frac{1-\rho}{\rho}}\;\frac{\varphi\big(u(x)\big)}{\varphi\big(\Phi^{-1}(x)\big)},\qquad
u(x)=\frac{\sqrt{1-\rho}\,\Phi^{-1}(x)-\Phi^{-1}(\PD)}{\sqrt{\rho}},
\label{eq:cc-density}
\end{equation}
which the calibration below maximises. With exposure $\EAD$ and an expected loss given default $\LGD$ independent of default, the loss
is $\EAD\cdot\LGD\cdot L_\infty$, the expected loss is $\EL=\EAD\cdot\LGD\cdot\PD$, and the risk measures at level $q$ are
\begin{equation}
\ES_q=\frac{\EAD\,\LGD}{1-q}\int_{-\infty}^{\Phi^{-1}(1-q)}p(z)\,\varphi(z)\dd z,\qquad
\mathrm{EC}_q=\VaR_q-\EL,
\label{eq:cc-es}
\end{equation}
where the expected shortfall integrates the conditional default probability over the worst $1-q$ of the factor, and the economic
capital $\mathrm{EC}_q$ is the unexpected loss, the same concept as the IRB capital requirement
\citep{credit-capital:basel2005,credit-capital:acerbitasche2002}.

\subsubsection{The IRB capital requirement}

For a corporate exposure the Basel Framework fixes the correlation $R$ and a maturity adjustment, and the capital requirement per unit
of exposure is
\begin{equation}
\begin{aligned}
K&=\left[\LGD\,\Phi\!\left(\frac{\Phi^{-1}(\PD)}{\sqrt{1-R}}+\sqrt{\frac{R}{1-R}}\,\Phi^{-1}(0.999)\right)-\PD\cdot\LGD\right]
\frac{1+(M - 2.5)\,b}{1-1.5\,b},\\
b&=\big(0.11852-0.05478\ln\PD\big)^2,
\end{aligned}
\label{eq:cc-irb}
\end{equation}
with $M$ the effective maturity in years and $\mathrm{RWA}=12.5\,K\cdot\EAD$ \citep{credit-capital:basel2005,credit-capital:basel-cre31}.
Writing the argument of $\Phi$ as $\big(\Phi^{-1}(\PD)+\sqrt{R}\,\Phi^{-1}(0.999)\big)/\sqrt{1-R}$ shows that the bracket is
$\LGD\,(\VaR_{0.999}-\PD)$ of \cref{eq:cc-lhp} with $\rho=R$: the IRB formula is the unexpected loss of the large homogeneous portfolio
at the regulatory correlation, times the maturity adjustment. The correlation depends on the class. With
$w(\PD)=(1-e^{-50\,\PD})/(1-e^{-50})$ and $w_{35}(\PD)=(1-e^{-35\,\PD})/(1-e^{-35})$,
\begin{equation}
\begin{aligned}
R_{\mathrm{corp}}&=0.12\,w+0.24\,(1-w),& R_{\mathrm{HVCRE}}&=0.12\,w+0.30\,(1-w),\\
R_{\mathrm{mort}}&=0.15,& R_{\mathrm{QRRE}}&=0.04,\\
R_{\mathrm{other}}&=0.03\,w_{35}+0.16\,(1-w_{35}),
\end{aligned}
\label{eq:cc-corr}
\end{equation}
for corporates, high-volatility commercial real estate, residential mortgages, qualifying revolving retail exposures and other retail;
the retail classes carry no maturity adjustment (CRE31.5 to CRE31.16; \citealp{credit-capital:basel-cre31}). The default probability
is floored by the regime (0.05 percent in the Basel III chapters, \citealp{credit-capital:basel-cre32}), and the regimes of the engine
differ in what they fix: Basel II multiplies the risk-weighted assets by 1.06 and floors the default probability at 0.03 percent
\citep{credit-capital:basel-cre30}, the European regulation applies the Basel III formula without the scaling \citep{credit-capital:crr3},
and the Chilean regulation uses the same risk function with correlations of its own and fixes the loss given default of the consumer
portfolio at 61.3 percent \citep{credit-capital:cmf-ran21-6}. \Cref{fig:cc-irb} draws the capital requirement of the five classes
against the default probability, and \cref{tab:cc-regimes} gives the risk weights under the four regimes at the fitted default
probabilities.

\begin{figure}[htbp]
\centering
\includegraphics[width=0.95\textwidth]{credit-capital-irb.pdf}
\caption{Capital requirement $K$ per unit of exposure at the 99.9 percent level against the probability of default, for the five IRB
classes of case C10 with an expected loss given default of 0.45 and, for the two non-retail classes, a maturity of 2.5 years (engine
\ccfn{regulatory.irb}). Filled markers are the IRB capital at each loan type's fitted default probability; open markers are the capital
of the large homogeneous portfolio at the fitted correlation, with the same maturity treatment.}
\label{fig:cc-irb}
\end{figure}

\begin{table}[htbp]
\centering
\small
\caption{Risk weight $12.5K$ in percent (columns Basel III, EU CRR3, Basel II and the Chilean RAN 21-6; the last column is Basel II over Basel III) for each loan type's IRB class at its fitted default probability, loss given default 0.45 and, for the non-retail classes, a maturity of 2.5 years, under the four regimes of the engine. Basel II carries the 1.06 scaling factor and a lower default probability floor; a dash marks a combination the engine refuses because the Chilean regime fixes the loss given default of the consumer portfolio at 61.3\%.}
\label{tab:cc-regimes}
\begin{tabular}{>{\raggedright\arraybackslash}p{2.8cm}>{\raggedright\arraybackslash}p{2.4cm}>{\raggedleft\arraybackslash}p{1.45cm}>{\raggedleft\arraybackslash}p{1.45cm}>{\raggedleft\arraybackslash}p{1.45cm}>{\raggedleft\arraybackslash}p{1.45cm}>{\raggedleft\arraybackslash}p{1.45cm}}
\toprule
Loan type & IRB class & Basel III & CRR3 & Basel II & Chile 21-6 & II / III \\
\midrule
Commercial and industrial & corporate & 89.2 & 89.2 & 94.6 & 89.2 & 1.060 \\
Commercial real estate & hvcre & 83.5 & 83.5 & 88.5 & 83.5 & 1.060 \\
Residential real estate & residential mortgage & 13.8 & 13.8 & 14.6 & 13.8 & 1.060 \\
Credit cards & qrre & 49.0 & 49.0 & 51.9 & -- & 1.060 \\
Other consumer & other retail & 45.3 & 45.3 & 48.0 & 45.3 & 1.060 \\
\bottomrule
\end{tabular}
\end{table}


\subsubsection{Calibration from yearly loss rates and the tail back-test}

A series of yearly loss rates $r_1,\dots,r_T$ is fitted for $(\PD,\rho)$ under the assumption that each year is an independent draw
$r_t=p(Z_t)$. The maximum-likelihood estimator maximises the log-likelihood of the density \cref{eq:cc-density},
\begin{equation}
\ell(\PD,\rho)=\sum_{t=1}^{T}\ln f(r_t;\PD,\rho),
\label{eq:cc-loglik}
\end{equation}
by the Nelder--Mead simplex on the logit scale of both parameters, with one restart from the optimum; a year without a rate is
excluded and counted, never interpolated, and at least three usable years are needed. The method-of-moments estimator uses the fact that
the second moment of the default fraction is the probability that two obligors default together,
\begin{equation}
\E[L_\infty^2]=\E\big[p(Z)^2\big]=\Phi_2\big(\Phi^{-1}(\PD),\Phi^{-1}(\PD);\rho\big)=\overline{r^2},\qquad \widehat{\PD}=\bar r,
\label{eq:cc-moments}
\end{equation}
with $\Phi_2$ the bivariate normal distribution function with correlation $\rho$. The joint default probability increases in $\rho$ from
$\PD^2$ to $\PD$, so the root is unique, and when $\overline{r^2}\le\PD^2$ there is no excess dispersion and the correlation is not
identified. The uncertainty of the maximum-likelihood estimate is the likelihood-ratio interval of the profile log-likelihood
$\ell_P(\rho)=\max_{\PD}\ell(\PD,\rho)$,
\begin{equation}
\Big\{\rho:\ \ell_P(\rho)\ge\ell_{\max}-\tfrac12\chi^2_{1,0.95}\Big\},\qquad \tfrac12\chi^2_{1,0.95}=1.920729,
\label{eq:cc-profile}
\end{equation}
and that of the moments estimate is a parametric-bootstrap interval of \ccBoot{} series of the same length drawn from the fitted model
\citep{credit-capital:vasicek2002}. \Cref{alg:cc-profile} states the procedure. The factor of a year implied by the fit is
\begin{equation}
\hat z_t=\frac{\Phi^{-1}(\PD)-\sqrt{1-\rho}\,\Phi^{-1}(r_t)}{\sqrt{\rho}},
\label{eq:cc-implied}
\end{equation}
the inverse of \cref{eq:cc-condpd}, decreasing in the loss rate. The fitted distribution is judged out of sample by counting the
evaluation years whose rate exceeds its fitted $q$-quantile: under the fit each year exceeds the $q$-quantile with probability $1-q$, so
the count is binomial and the upper tail probability of the observed count $x_{\mathrm{obs}}$ out of $n$ years is
\begin{equation}
p_{\mathrm{back}}=\Pr\big(X\ge x_{\mathrm{obs}}\big),\qquad X\sim\mathrm{Binomial}(n,\,1-q).
\label{eq:cc-backtest}
\end{equation}
The test has little power with $n=\ccEvYears$ years and the $p$-value is therefore shown beside the count.

\begin{algorithm}[htbp]
\caption{Maximum-likelihood fit of $(\PD,\rho)$ with the profile-likelihood interval of the asset correlation.}
\label{alg:cc-profile}
\begin{algorithmic}[1]
\Require yearly loss rates $r_1,\dots,r_T$ (missing years removed, $T\ge3$), level $1-\alpha$
\State $(\theta_1,\theta_2)\gets$ Nelder--Mead on $-\ell$ over $(\operatorname{logit}\PD,\operatorname{logit}\rho)$, started at $(\operatorname{logit}\bar r,\operatorname{logit}0.1)$
\State restart Nelder--Mead from the optimum; $(\widehat{\PD},\widehat\rho)\gets(\operatorname{expit}\theta_1,\operatorname{expit}\theta_2)$
\State $c\gets\ell(\widehat{\PD},\widehat\rho)-\tfrac12\chi^2_{1,1-\alpha}$
\For{each side of $\widehat\rho$ (toward $10^{-4}$ and toward $0.9999$)}
\State step along a grid of 41 points on the logit scale
\State \textbf{at each point} $\rho$: maximise $\ell(\PD,\rho)$ over $\PD\in(10^{-6},0.99)$ (grid on the logit scale, then bounded refinement)
\State stop at the first grid cell where $\ell_P(\rho)<c$ and solve $\ell_P(\rho)=c$ inside it by Brent's method
\State if the grid ends first, return the bound of the parameter space
\EndFor
\State \Return $\widehat\rho$ and the two endpoints (the interval of $\PD$ is obtained the same way)
\end{algorithmic}
\end{algorithm}

\paragraph{Assumptions.} The model assumes one Gaussian systematic factor; an infinitely granular and homogeneous portfolio; a loss given
default that is constant and independent of default; losses from default only, over a horizon of one year; and yearly loss rates that are
independent draws of $p(Z)$ from a stationary distribution. The charge-off rate stands for the portfolio's yearly loss rate, which it
only approximates: it is net of recoveries and is a rate of charge-offs, not of defaults. The IRB class is the nearest class of the
Basel Framework and not a claim that the Board's series is that class; the Board's commercial real estate series is not the
high-volatility class, which is a stricter subset.

\paragraph{Parameters in the case.} The confidence levels are 99 and 99.9 percent; the expected loss given default is \ccLgd{} for every
name, a choice that isolates the effect of correlation, concentration and tail dependence; the maturity of the non-retail classes is 2.5
years, the supervisory default; the regime is the Basel III chapter, with the default probability checked against its floor so that the
comparison uses the fitted value unchanged. The fit window is \ccFitFirst{} to \ccFitLast{} (\ccFitYears{} years, or the years
\cref{tab:cc-cases} lists where the series starts later), the evaluation window \ccEvFirst{} to \ccEvLast{}, the bootstrap has \ccBoot{}
draws, and the profile is evaluated on 41 correlations around the estimate plus the estimate itself. A fit that is handed an
evaluation-window year is refused, which keeps the tail test out of sample.

\paragraph{Complexity.} The distribution function, quantile and density are closed forms, $O(1)$ each; the expected shortfall is a
Simpson quadrature of 4,000 steps on $[\Phi^{-1}(1-q)-12,\,\Phi^{-1}(1-q)]$. One likelihood evaluation costs $O(T)$, so a
Nelder--Mead run costs $O(T)$ times its evaluation count (at most 8,000), a profile point costs $O(T)$ times the points of its inner
search, and the bootstrap costs $O(B)$ fits of the moments estimator, each a root finding in which every step is one quadrature. The IRB
requirement is $O(1)$ per exposure, $O(N)$ for $N$ exposures.

\paragraph{Implementation.} The conditional default probability, distribution function, density, quantile and implied factor are in
\ccfn{riskvalidation.engines.vasicek}, the loss measures of the large homogeneous portfolio in
\ccfn{riskvalidation.portfolio.portfolio.lhp_measures}, the estimators and intervals in \ccfn{riskvalidation.portfolio.asrf_fit}
(\ccfn{fit_ml}, \ccfn{fit_moments}, \ccfn{profile_interval}, \ccfn{bootstrap_interval}, \ccfn{implied_factors}), the
tail back-test in \ccfn{riskvalidation.validation.portfolio.tail_backtest} and the IRB requirement in
\ccfn{riskvalidation.regulatory.irb.capital_requirement}, all in release \ccEngine{} of the engine. The case code
(\ccfn{data-pipeline/pipeline/cases/c10/documents.py}) orders the engine's results and computes none of them.

\paragraph{Comparison and acceptance.} The large-portfolio capital at the fitted correlation is compared with the IRB capital of the
nearest class at the same default probability and with the same maturity treatment, so that their ratio isolates the correlation
(\cref{tab:cc-irb}). The acceptance criteria are those of the published tables. The engine reproduces the Basel Committee's illustrative
risk weights of CRE99 (eight exposure types at 18 default probabilities, 144 cells; \citealp{credit-capital:basel-cre99}): 141 cells agree
with the printed two decimals and three differ by one unit in the last place, so the test applies a tolerance of one printed unit, 0.01
percentage points, which is looser than the half-unit (0.005 percentage points) that an exact reproduction would meet. The Basel II regime
carries the 1.06 scaling and the 0.03 percent floor and is tested to equal 1.06 times the Basel III value on every cell, the European
regime to equal it, and the Chilean regime to use the correlations of its numeral 4.4.2. The loss quantities of the large portfolio
reproduce closed forms to the printed digits (\cref{tab:cc-anchors}): for a default probability of 0.02 and a correlation of 0.15 the
99.9 percent value at risk is \ccAnchorVar{}, the expected shortfall \ccAnchorEs{}, and $0.45\,(\ccAnchorVar{}-0.02)=\ccAnchorK{}$ is the
IRB capital of the mortgage class.

\begin{table}[htbp]
\centering
\small
\caption{Analytic anchors of the portfolio engine recomputed by the table script: the engine value, an independent closed form or numerical integral computed with scipy (or the value printed in the cited paper), and the absolute difference. The first three rows use a default probability of 0.02, a correlation of 0.15 and, for the loss quantities, unit exposure and unit loss; the last four rows use the small portfolios stated in the row.}
\label{tab:cc-anchors}
\begin{tabular}{>{\raggedright\arraybackslash}p{7.5cm}>{\raggedleft\arraybackslash}p{1.9cm}>{\raggedleft\arraybackslash}p{2.1cm}>{\raggedleft\arraybackslash}p{2.3cm}}
\toprule
Quantity & Engine & Independent value & Absolute difference \\
\midrule
\relax LHP value at risk, 99.9\% & 0.176329 & 0.176329 & \ensuremath{0} \\
\relax LHP expected shortfall, 99.9\% (Simpson, 4,000 steps) & 0.209404 & 0.209404 & \ensuremath{1.8\times10^{-12}} \\
\relax IRB capital K (mortgage class) against LGD $\times$ (LHP VaR $-$ PD) & 0.070348 & 0.070348 & \ensuremath{1.4\times10^{-17}} \\
\relax CreditRisk+ recursion, one obligor, negative binomial (max over 4 terms) & 0.826446 & 0.826446 & \ensuremath{2.2\times10^{-16}} \\
\relax CreditRisk+ recursion, three obligors without sector, against the convolution of Poisson laws (max over 8 terms) & 0.548812 & 0.548812 & \ensuremath{6.9\times10^{-18}} \\
\relax Granularity multiplier $\delta$ at $\xi=0.25$, 99.9\% against the value printed by Gordy and L\"utkebohmert & 4.8336 & 4.83 & 0.0036 \\
\relax Exponential twist of a default probability 0.1 at $\theta=\ln 2$ & 0.181818 & 0.181818 & \ensuremath{5.6\times10^{-17}} \\
\bottomrule
\end{tabular}
\end{table}


\subsection{Portfolio loss models, measures, granularity and allocation}
\label{sec:m17}

\subsubsection{The portfolio, its loss and its risk measures}

The portfolio loss models share one object and one set of measures. Over a one-year horizon the default-mode loss of obligors
$i=1,\dots,N$ with default indicators $D_i$ is
\begin{equation}
L=\sum_{i=1}^{N}\EAD_i\,\LGD_i\,D_i,\qquad \EL=\E[L]=\sum_i\EAD_i\,\E[\LGD_i]\,\PD_i,
\label{eq:cc-loss}
\end{equation}
the second form holding when the loss given default is independent of default \citep{credit-capital:acerbitasche2002}. The measures are the
value at risk $\VaR_q(L)=\inf\{\ell:\Pr(L\le\ell)\ge q\}$, the expected shortfall of \cref{eq:cc-es} (equal to $\E[L\mid L\ge\VaR_q]$ for a
continuous loss) and the economic capital $\VaR_q-\EL$. For a sample of $S$ losses sorted ascending, $L_{(1)}\le\dots\le L_{(S)}$, and
$k=\lceil qS - 10^{-9}\rceil$, the estimators are
\begin{equation}
\widehat{\VaR}=L_{(k)},\qquad
\widehat{\ES}=\frac{1}{1-q}\left[\sum_{i>k}\frac{L_{(i)}}{S}+\Big(\frac kS-q\Big)L_{(k)}\right],
\label{eq:cc-sample-es}
\end{equation}
in which the boundary scenario carries the fractional probability mass so that the tail has mass exactly $1-q$ (the weights of
\citealp{credit-capital:acerbitasche2002}). The Monte Carlo error of both is the batch-means standard error: the sample is cut into 20
consecutive equal batches, both estimators are computed in each, and the error is the standard deviation of the batch estimates divided by
$\sqrt{20}$. A value at risk on a lattice (a portfolio whose loss is a whole number of equal name losses) can be identical in every batch,
and its standard error is then exactly zero; that is a lattice and not a certain quantile.

\subsubsection{Copula simulation: Gaussian, multi-factor Gaussian and Student-t}

The Gaussian copula generalises \cref{eq:cc-asset} to $F$ systematic factors, as in \citet{credit-capital:li2000}. With $Z\sim N(0,\Sigma)$
($\Sigma$ the correlation matrix of the factors), loadings $\beta_i$ and independent standard normal $e_i$, obligor $i$ defaults when
\begin{equation}
A_i=\beta_i^{\top}Z+\sqrt{1-\beta_i^{\top}\Sigma\,\beta_i}\;e_i<\Phi^{-1}(\PD_i),
\label{eq:cc-gaussian}
\end{equation}
and loses $\EAD_i\LGD_i$; the idiosyncratic scale uses $\beta_i^{\top}\Sigma\beta_i$ so that $A_i$ has unit variance, which equals
$\beta_i^\top\beta_i$ for one factor and for uncorrelated factors. Gordy compares the structural family of such copula models and the actuarial family of
CreditRisk+ in one framework \citep{credit-capital:gordy2000}. Migration-mode models such as CreditMetrics revalue each exposure at the horizon on
the curve of its new rating \citep{credit-capital:creditmetrics1997}; the engine implements the default mode only. The Student-t copula multiplies the
asset return by a common mixing variable,
\begin{equation}
A_i=\sqrt{W}\Big(\beta_i^{\top}Z+\sqrt{1-\beta_i^{\top}\Sigma\,\beta_i}\;e_i\Big)<t_\nu^{-1}(\PD_i),\qquad W=\nu/\chi^2_\nu,\ \nu>2,
\label{eq:cc-student}
\end{equation}
with $W$ independent of $Z$ and $e$ and $t_\nu^{-1}$ the quantile function of the Student distribution with $\nu$ degrees of freedom. It
adds tail dependence: joint defaults cluster more than under the Gaussian copula for the same asset correlation, and the marginals are
unchanged \citep{credit-capital:freymcneil2003,credit-capital:demartamcneil2005}. For two names with default probability 0.01 and
correlation 0.3 the Gaussian joint default probability is 0.000556 and the Student-t value with $\nu=5$ is 0.001647. \Cref{alg:cc-copula}
states the simulation.

\begin{algorithm}[htbp]
\caption{Chunked copula Monte Carlo of the portfolio loss with Euler tracking by group.}
\label{alg:cc-copula}
\begin{algorithmic}[1]
\Require portfolio $(\EAD,\PD,\LGD,\beta,\Sigma)$, scenario count $S$, seed, chunk size $c$, optional $\nu$
\State draw from $\mathrm{PCG64}(\text{seed})$; $C\gets\operatorname{chol}(\Sigma)^{\top}$; threshold $\theta_i\gets\Phi^{-1}(\PD_i)$ (or $t_\nu^{-1}(\PD_i)$)
\For{each block of $\lfloor c/3\rfloor$ scenarios}
\State $Z\gets$ standard normals $\times\,C$; $A\gets$ standard normals $\times$ idiosyncratic scale $+\,Z\beta^{\top}$
\State if Student-t: $A\gets A\sqrt{\nu/\chi^2_\nu}$ with one $\chi^2_\nu$ draw per scenario
\State $D\gets\mathbf{1}\{A<\theta\}$; scenario loss $\gets D\cdot(\EAD\cdot\LGD)$; group loss $\gets$ scenario loss summed within each group
\EndFor
\State sort the $S$ losses; \Return $\widehat{\VaR},\widehat{\ES}$ of \cref{eq:cc-sample-es} and the batch-means errors
\end{algorithmic}
\end{algorithm}

The result is reproducible for the same seed, scenario count and chunk size; a different chunk size gives a different, equally valid
sample. The block size bounds the working memory at $c\,N$ double-precision numbers, which the case sets to at most 150 million bytes per
chunk. The time is $O(S\,N\,F)$ for the matrix products and $O(S\log S)$ for the sort, and the memory is $O(S)$ for the losses plus
$O(SG)$ for the losses of $G$ groups. The Monte Carlo standard error falls as $1/\sqrt S$.

\subsubsection{CreditRisk+ and its stable recursion}

CreditRisk+ models the number of defaults as Poisson with a stochastic intensity. Exposures are integer multiples $\nu_i$ of a loss unit,
$X=\sum_i\nu_iY_i$ with $Y_i$ Poisson of intensity $p_i\big(w_{0,i}+\sum_{k=1}^{K}w_{k,i}S_k\big)$, where $w_{0,i}+\sum_kw_{k,i}=1$ and
the sector factors $S_k$ are independent gamma variables with mean one and variance $\sigma_k^2$. The probability generating function of
the loss in units is \citep{credit-capital:haaf2004}
\begin{equation}
G(z)=\exp\!\Bigg(\sum_{i}w_{0,i}\,p_i\,(z^{\nu_i}-1)-\sum_{k=1}^{K}\frac{1}{\sigma_k^2}\ln\Big[1-\sigma_k^2\sum_{i}w_{k,i}\,p_i\,(z^{\nu_i}-1)\Big]\Bigg).
\label{eq:cc-pgf}
\end{equation}
Write $P_k(z)=\sum_iw_{k,i}p_iz^{\nu_i}$, $\mu_k=P_k(1)$, $Q_k=P_k/\mu_k$ and $t_k=\sigma_k^2\mu_k/(1+\sigma_k^2\mu_k)$. Then
$1-\sigma_k^2(P_k-\mu_k)=(1+\sigma_k^2\mu_k)(1-t_kQ_k)$ and $-\ln(1-t_kQ_k)=\sum_{j\ge1}b_{kj}z^j$, so that $\ln G(z)=\sum_ja_jz^j$ with
\begin{equation}
\begin{aligned}
a_0&=-\sum_iw_{0,i}p_i-\sum_k\frac{\ln(1+\sigma_k^2\mu_k)}{\sigma_k^2},\qquad
a_j=\sum_{i:\nu_i=j}w_{0,i}p_i+\sum_k\frac{b_{kj}}{\sigma_k^2},\\
b_{kj}&=t_k\,q_{kj}+\frac{t_k}{j}\sum_{l<j}(j-l)\,b_{k,j-l}\,q_{kl},
\end{aligned}
\label{eq:cc-crp-coefficients}
\end{equation}
where $q_{kj}$ is the coefficient of $z^j$ in $Q_k$; the recursion for $b_{kj}$ comes from the derivative of the series. From
$G'=(\ln G)'G$ the probabilities $p_n=\Pr(X=n)$ follow:
\begin{equation}
p_0=e^{a_0},\qquad n\,p_n=\sum_{j=1}^{n}j\,a_j\,p_{n-j}.
\label{eq:cc-crp-recursion}
\end{equation}
Every term is nonnegative, so no subtraction of terms of opposite sign occurs and no probability is negative whatever the grid length
\citep{credit-capital:haaf2004}. The Panjer recursion advocated in the original CreditRisk+ document clears the denominators and subtracts
such terms of opposite sign, and \citet{credit-capital:haaf2004} cite it as numerically unstable. The engine keeps it
only as a comparison: on a portfolio of 300 obligors in 40 exposure bands with two sectors it agrees with the stable recursion to a relative
$3\times10^{-14}$ wherever the stable probability exceeds $10^{-250}$, and on four larger random portfolios the relative difference reached
$10^{-11}$ with no negative entry, so the reported instability is not reproduced and the stable recursion is the only result used.
\Cref{alg:cc-crp} states the computation.

\begin{algorithm}[htbp]
\caption{Stable CreditRisk+ recursion for the loss distribution on an integer grid of loss units.}
\label{alg:cc-crp}
\begin{algorithmic}[1]
\Require exposures $\nu_i$ in loss units, default probabilities $p_i$, weights $w_{0,i},w_{k,i}$, variances $\sigma_k^2$, grid length $n_{\max}$
\State bin the obligors by exposure band: coefficients of $\sum_i w_{k,i}p_i z^{\nu_i}$ for $k=0,\dots,K$
\State $a_0\gets$ \cref{eq:cc-crp-coefficients}; $a_j\gets\sum_{i:\nu_i=j}w_{0,i}p_i$ for $j\ge1$
\For{each sector $k$ with positive weight}
\State $\mu_k,\ q_{k\cdot},\ t_k$ as defined; $b_{k0}\gets0$
\For{$j=1,\dots,n_{\max}$}
\State $b_{kj}\gets t_kq_{kj}+(t_k/j)\sum_{l<j}(j-l)\,b_{k,j-l}\,q_{kl}$ over the occupied exposure bands $l$
\EndFor
\State $a_j\gets a_j+b_{kj}/\sigma_k^2$ for all $j$
\EndFor
\State $p_0\gets e^{a_0}$ (raise if $e^{a_0}<10^{-300}$); \textbf{for} $n=1,\dots,n_{\max}$ \textbf{do} $p_n\gets\frac1n\sum_{j=1}^{n}j\,a_j\,p_{n-j}$
\State \Return $p_0,\dots,p_{n_{\max}}$ (the mass beyond the grid is reported, never added to the last entry)
\end{algorithmic}
\end{algorithm}

The cost is $O\big(K\,(N+n_{\max}B)\big)$ to bin the exposures and compute the coefficients $b_{kj}$, with $B$ the number of occupied
exposure bands, and $O(n_{\max}^2)$ for the final recursion, with $O(n_{\max})$ memory. In the case there is one sector ($K=1$, $w_0=0$) whose variance matches the variance of
the Vasicek default rate, $\sigma^2=\big(\E[p(Z)^2]-\PD^2\big)/\PD^2$, computed by Gauss--Hermite quadrature of the engine's conditional
default probability, and the loss unit is \ccLossUnit{} USD; the grid doubles until the mass beyond it is below $10^{-7}$.

\subsubsection{The granularity adjustment}

The large-portfolio quantile ignores that a finite portfolio has idiosyncratic risk left. Let $L_n=\sum_is_iU_i$ be the loss rate with
exposure shares $s_i$, let $X$ be the systematic factor with density $h$, and $\mu(x)=\E[L_n\mid X=x]$, $\sigma^2(x)=\mathrm V[L_n\mid X=x]$.
\citet{credit-capital:gordylutkebohmert2013} (their eq. 9) show that the gap of the portfolio quantile above the infinitely granular one,
already a share of the exposure, is
\begin{equation}
\mathrm{GA}=\alpha_q(L_n)-\mu\big(\alpha_q(X)\big)\approx-\frac{1}{2h(\alpha_q(X))}\left.\frac{\mathrm d}{\mathrm dx}\left[\frac{h(x)\,\sigma^2(x)}{\mu'(x)}\right]\right|_{x=\alpha_q(X)}.
\label{eq:cc-ga-general}
\end{equation}
The Basel form is built on the single-factor CreditRisk+ model with a gamma factor of mean one and variance $1/\xi$, conditional default
probability $\pi_i(x)=\PD_i(1+\omega_i(x-1))$, reserve $R_i=\E[\LGD_i]\PD_i$ and the capital $K_i$ of name $i$; with
$K^{*}=\sum_is_iK_i$, $C_i=(\mathrm V[\LGD_i]+\E[\LGD_i]^2)/\E[\LGD_i]$, $\mathrm V[\LGD_i]=\gamma\,\E[\LGD_i](1-\E[\LGD_i])$ and
$\alpha_q=\alpha_q(X)$ the $q$-quantile of the gamma factor, the simplified adjustment (their eqs. 11, 15, 16 and 17) is
\begin{equation}
\widetilde{\mathrm{GA}}=\frac{1}{2K^{*}}\sum_{i=1}^{n}s_i^{2}\,C_i\big(\delta\,(K_i+R_i)-K_i\big),\qquad
\delta=(\alpha_q-1)\Big(\xi+\frac{1-\xi}{\alpha_q}\Big).
\label{eq:cc-ga}
\end{equation}
The paper's test vector, $\xi=\ccXi$ and $q=0.999$, gives $\alpha_q=\ccAlphaQ$ and $\delta=\ccDelta$ (the paper prints $\ccDeltaPaper$), and
the other reading of $\xi$ (a factor of variance $\xi$) gives 1.0867. The capital $K_i$ is an input, as in the regulatory algorithm of the
paper: the case takes it from the IRB formula at the regulatory correlation, which eliminates the default probability and the weight
$\omega_i$ from the inputs. The exact adjustment (their eq. 14), with second-order terms, and the extension to mark-to-market models
\citep{credit-capital:gordymarrone2012} are not implemented, and the adjustment is zero in the limit of many names, being proportional to $1/n$ for $n$ identical names.

The reference for the adjustment is the CreditRisk+ recursion of the same single-factor model for $n$ identical names with the weight
$\omega=K/\big(\LGD\,\PD\,(\alpha_q-1)\big)$ implied by the IRB capital of one name \citep{credit-capital:gordylutkebohmert2013}; it is
defined only where $\omega\in(0,1]$. Its recoveries are deterministic, so the reference leaves out the term of the formula that carries the
variance of the loss given default.

\subsubsection{Importance sampling of the Gaussian copula}

Plain Monte Carlo at the 99.9 percent level has about $10^{-3}S$ scenarios in the tail. \citet{credit-capital:glassermanli2005} sample the
Gaussian copula under a measure that makes the tail common. Given the factor value $z$, the conditional default probabilities
$p_i(z)$ are twisted, with $c_i=\EAD_i\LGD_i$, to
\begin{equation}
p_{i,\theta}(z)=\frac{p_i(z)\,e^{\theta c_i}}{1+p_i(z)\,(e^{\theta c_i}-1)},
\label{eq:cc-twist}
\end{equation}
with $\theta=\theta(z)\ge0$ the root of $\sum_ic_ip_{i,\theta}(z)=x$ for the threshold $x$ ($\theta=0$ where the untwisted mean loss is at or
above $x$). The factor is drawn with a shifted mean $\mu$, and each scenario carries the likelihood ratio
\begin{equation}
w=\exp\Big(-\theta L+\sum_i\ln\big(1+p_i\,(e^{\theta c_i}-1)\big)\Big)\exp\Big(-\mu^{\top}z+\tfrac12\mu^{\top}\mu\Big),
\label{eq:cc-weight}
\end{equation}
which is exact whatever $\theta(z)\ge0$ and $\mu$ are, so that $\E_Q[w]=1$ and $\E_Q[wL]=\EL$. The default shift maximises
$-\theta(z)\,x+\ln M(\theta(z),z)-z^{\top}z/2$ ($M$ the moment generating function of the twisted loss) by Nelder--Mead from $z=0$, a
deterministic search. The transcription of the twist and of the mean shift was not checked against the paper, which is behind a paywall,
so the engine accepts the estimator by one criterion: it must be unbiased against plain Monte Carlo (the weights average one, the weighted
loss averages the expected loss, and the weighted tail probability, value at risk and expected shortfall agree with the plain sample within
four combined standard errors). The weighted quantile is the smallest loss whose weighted upper tail $\sum_iw_i\mathbf 1\{L_i>x\}/S$ is at
most $1-q$, normalised by the scenario count $S$ and never by the realised sum of the weights, which is far from $S$ for a threshold deep in
the tail. A heavy-tailed variant exists for risk factors with heavy tails \citep{credit-capital:ghs2002}. \Cref{alg:cc-is} states the procedure. Only the Gaussian copula with uncorrelated factors is supported; there is no Student-t
twist, so that rung is reported as not evaluated and no number stands in for it.

\begin{algorithm}[htbp]
\caption{Importance sampling of the Gaussian copula loss by the exponential twist and a shifted factor.}
\label{alg:cc-is}
\begin{algorithmic}[1]
\Require portfolio with uncorrelated factors, threshold $x$ (the plain Monte Carlo 99.9 percent value at risk), scenario count $S$, seed
\State $\mu\gets\arg\max_z\big[-\theta(z)x+\ln M(\theta(z),z)-z^{\top}z/2\big]$ by Nelder--Mead from $0$
\For{each block of scenarios}
\State $z\gets\mu+$ standard normal; $p_i\gets\Phi\big((\Phi^{-1}(\PD_i)-\beta_i^{\top}z)/\sqrt{1-\beta_i^{\top}\beta_i}\big)$
\State $\theta\gets$ safeguarded Newton on $\sum_ic_ip_{i,\theta}=x$ for the rows whose mean loss is below $x$, else $0$
\State draw defaults with the twisted $p_{i,\theta}$ of \cref{eq:cc-twist}; $L\gets\sum_ic_iD_i$; $\log w\gets$ \cref{eq:cc-weight}
\EndFor
\State \Return $\widehat{\VaR},\widehat{\ES}$ with the weights $w$, normalised by $S$, and their batch-means errors
\end{algorithmic}
\end{algorithm}

\subsubsection{Euler allocation and concentration}

For a positively homogeneous and differentiable risk measure $\varrho$ and the loss $L(u)=\sum_iu_iL_i$, the Euler principle allocates
the capital in full \citep{credit-capital:tasche2008}:
\begin{equation}
\begin{aligned}
\varrho(L)&=\sum_iu_i\left.\frac{\partial\varrho(L(u))}{\partial u_i}\right|_{u=1},\\
\mathrm{RC}_i^{\VaR}&=\E\big[L_i\mid L=\VaR_q(L)\big],\qquad
\mathrm{RC}_i^{\ES}=\E\big[L_i\mid L\ge\VaR_q(L)\big].
\end{aligned}
\label{eq:cc-euler}
\end{equation}
Conditioning on a single loss level is very noisy in simulation, so the expected-shortfall contribution is computed with the weights
of \cref{eq:cc-sample-es}, $\mathrm{RC}_g=\frac1{1-q}\sum_sw_sL_{g,s}$, which adds up to the estimated expected shortfall exactly, and the
value-at-risk contribution with the Nadaraya--Watson kernel estimator of \citet{credit-capital:tasche2009},
\begin{equation}
\widehat{\mathrm{RC}}_g=\frac{\sum_s\kappa\big((L_s-\widehat{\VaR})/h\big)L_{g,s}}{\sum_s\kappa\big((L_s-\widehat{\VaR})/h\big)},\qquad h=1.06\,\mathrm{sd}(L)\,S^{-1/5},
\label{eq:cc-kernel}
\end{equation}
with $\kappa$ the Gaussian kernel and a bandwidth of Silverman's rule (the contract's choice, since the paper names none); the kernel sum
agrees with the sample value at risk within a stated tolerance and not exactly. The asymptotic anchor is the first term of
\citet{credit-capital:emmertasche2005}: the loss quantile of the one-factor model is the conditional mean at $z=\Phi^{-1}(1-q)$, so
$\mathrm{RC}_i=\EAD_i\LGD_i\,p_i\big(\Phi^{-1}(1-q)\big)$ with $p_i$ of \cref{eq:cc-condpd}. It sums to the infinitely granular value at
risk and is portfolio-invariant; it is not their full contribution, which adds granularity terms. Name concentration is summarised by the
Herfindahl index $\mathrm{HHI}=\sum_is_i^2$ and the effective number of names $1/\mathrm{HHI}$; the granularity adjustment is driven by an
$s_i^2$-weighted sum, so it is a model-based refinement of the index. Sector concentration \citep{credit-capital:dullmannmasschelein2007}
is not implemented. Euler allocation costs $O(S\,G)$ plus the sort for the groups of the sample.

\subsubsection{Validation tests of the portfolio models}

Five tests judge the outputs, each with a null hypothesis and a light. The tail back-test of \cref{eq:cc-backtest} has the null that the
count of evaluation years above the fitted quantile is binomial. The Monte Carlo convergence test has the null that the Monte Carlo
estimate equals the analytic anchor, with statistic $z=(\text{estimate}-\text{anchor})/\mathrm{se}$ and two-sided $p=2(1-\Phi(|z|))$; the
anchor is exact only for the large homogeneous portfolio, so the case runs it on 5,000 identical names and the granularity gap of the
finite portfolio is the adjustment's, not the test's. The Euler additivity test has the metric $|\sum\mathrm{RC}-\text{total}|/|\text{total}|$,
green below a tolerance. The recursion test compares the CreditRisk+ distribution with a Monte Carlo sample of the same Poisson--gamma
model by the Kolmogorov distance $D=\max_n|F_{\mathrm{pmf}}(n)-F_{\mathrm{mc}}(n)|$ and its asymptotic law
$p=2\sum_{j\ge1}(-1)^{j-1}e^{-2j^2SD^2}$, which is conservative for a discrete null; any negative probability makes it red. The size and
power test gives the rejection rate, its standard error $\sqrt{r(1-r)/R}$ over $R$ replications and the Wilson interval
\citep{credit-capital:wilson1927}. The lights are green when the $p$-value is at least 0.05, amber between 0.01 and 0.05 and red below.
These are acceptance criteria of the validation, not regulatory thresholds, and a green light says only that the evidence does not
reject the null.

\paragraph{Assumptions.} Default-mode losses over one year; no migration loss and no random loss given default in the simulations; Gaussian
or Student-t dependence with the correlation of the fit; for the Student-t copula the degrees of freedom are an input and are poorly
identified from data; for the sector rung the correlation between the two sectors is an input.

\paragraph{Parameters in the case.} Every model uses the fitted default probability and correlation of its loan type. The copula samples
have \ccScen{} scenarios with seeds \ccSeedOne{}, \ccSeedMf{} and \ccSeedStu{} (one factor, two sectors, Student-t); the two-sector rung has
a correlation of \ccSectorCorr{} between its sector factors; the Student-t copula has $\nu=\ccNu$; importance sampling has \ccScenIs{}
scenarios with seed \ccSeedIs{} and its threshold is the plain Monte Carlo 99.9 percent value at risk; CreditRisk+ has the loss unit of
\ccLossUnit{} USD; the granularity adjustment has $\xi=\ccXi$, $\gamma=\ccGamma$ and capital $K_i$ from the IRB formula at the group's default
probability; the expected loss given default is \ccLgd{}. The illustrative portfolio is \ccTotalEad{} million USD in three groups of identical
names (large, mid-size and small), with group default probabilities $0.5$, $1$ and $1.5$ times the fitted one, rescaled so that the
exposure-weighted default probability is exactly the fitted one, and every name loading on the factor with $\sqrt{\hat\rho}$; the names and
exposures per loan type are in \cref{tab:cc-cases}. The values were fixed to isolate the effects the case studies: one loss given default
removes its variation, three groups give an allocation to read, and a mix of few large names for commercial loans and many small names for
consumer loans gives the range of concentration (\ccEffNamesCre{} to \ccEffNamesCard{} effective names).

\paragraph{Implementation.} The portfolio record and the large-portfolio measures are in \ccfn{riskvalidation.portfolio.portfolio}, the
copulas in \ccfn{riskvalidation.portfolio.copula} (\ccfn{simulate_gaussian}, \ccfn{simulate_student_t}, \ccfn{mc_standard_errors}),
CreditRisk+ in \ccfn{riskvalidation.portfolio.creditrisk_plus} (\ccfn{pmf_stable}, \ccfn{pmf_panjer}, \ccfn{measures_from_pmf}),
the adjustment in \ccfn{riskvalidation.portfolio.granularity}, importance sampling in \ccfn{riskvalidation.portfolio.importance}, the
allocation and concentration in \ccfn{riskvalidation.portfolio.euler} and the tests in \ccfn{riskvalidation.validation.portfolio}, in release
\ccEngine{} of the engine. The case orders the engine's answers in \ccfn{data-pipeline/pipeline/cases/c10/ladder.py}; every loss measure,
contribution, adjustment and test row is returned by the engine.

\paragraph{Comparison and acceptance.} The models are compared with each other on the same portfolio, and the finite-portfolio models with
the large-portfolio anchor. The acceptance criteria are: the Monte Carlo converges to the large-portfolio quantile as the number of names
grows; the Euler contributions sum to the total (exactly for the expected shortfall, within the stated 1 percent tolerance for the kernel
value at risk); the CreditRisk+ recursion matches a Monte Carlo of the same model; the multiplier $\delta$ reproduces the paper's 4.83;
and the importance-sampling estimates are unbiased against plain Monte Carlo. The anchors the engine reproduces are in
\cref{tab:cc-anchors}, and the results of the tests on the five portfolios are in \cref{tab:cc-battery}.

% results
\subsection{Results}
\label{sec:results-credit-capital}

The results follow the order of the methods: the fit and the correlation, the out-of-sample tail, the comparison of the models on the
illustrative portfolios, the granularity adjustment and the allocation, the validation tests and, last, the size and power of the
correlation estimators on the simulator.

\subsubsection{The fitted correlation and the IRB capital}

\begin{table}[htbp]
\centering
\small
\caption{Vasicek fit of the yearly net charge-off rates on the fit window: the maximum-likelihood estimate of the default probability (percent of the balance per year) and of the asset correlation with their 95\% profile-likelihood intervals, and the method-of-moments correlation with its 95\% parametric-bootstrap interval of 500 draws; intervals in square brackets.}
\label{tab:cc-fit}
\begin{tabular}{>{\raggedright\arraybackslash}p{2.6cm}>{\raggedleft\arraybackslash}p{0.9cm}>{\raggedright\arraybackslash}p{3.0cm}>{\raggedright\arraybackslash}p{3.3cm}>{\raggedright\arraybackslash}p{3.3cm}}
\toprule
Loan type & Years & PD, ML (\%) & $\rho$, ML & $\rho$, moments \\
\midrule
Commercial and industrial & 22 & $0.910$ [$0.680$, $1.297$] & $0.064$ [$0.038$, $0.116$] & $0.043$ [$0.019$, $0.076$] \\
Commercial real estate & 16 & $0.436$ [$0.188$, $1.336$] & $0.238$ [$0.144$, $0.407$] & $0.156$ [$0.052$, $0.229$] \\
Residential real estate & 16 & $0.139$ [$0.116$, $0.173$] & $0.013$ [$0.007$, $0.028$] & $0.013$ [$0.005$, $0.023$] \\
Credit cards & 22 & $4.228$ [$3.815$, $4.717$] & $0.013$ [$0.007$, $0.024$] & $0.012$ [$0.006$, $0.022$] \\
Other consumer & 22 & $0.975$ [$0.866$, $1.110$] & $0.011$ [$0.006$, $0.021$] & $0.010$ [$0.005$, $0.018$] \\
\bottomrule
\end{tabular}
\end{table}


\Cref{tab:cc-fit} gives the maximum-likelihood fit with its profile interval and the moments fit with its bootstrap interval. The fitted
asset correlations are low for four of the five series: \ccRhoMlRes{} for residential real estate (95 percent interval \ccRhoMlLoRes{} to
\ccRhoMlHiRes{}), \ccRhoMlCard{} for credit cards, \ccRhoMlOth{} for other consumer and \ccRhoMlCi{} for commercial and industrial loans; only
commercial real estate gives \ccRhoMlCre{} (\ccRhoMlLoCre{} to \ccRhoMlHiCre{}), on sixteen years that include the very low rates of 1997 and 1998. The intervals are wide, with the ratio of the upper to the lower bound between \ccRhoWidthMin{} and \ccRhoWidthMax{},
which is what 16 to \ccFitYears{} yearly points can identify; \cref{fig:cc-profile} shows the profile log-likelihoods on which the intervals
rest. The method of moments agrees with the maximum-likelihood estimate for the consumer and residential series (\ccRhoMoRes{} against
\ccRhoMlRes{} for residential real estate, \ccRhoMoCard{} against \ccRhoMlCard{} for credit cards) and gives smaller values for the two
commercial series (\ccRhoMoCi{} against \ccRhoMlCi{} and \ccRhoMoCre{} against \ccRhoMlCre{}); the intervals of the two estimators overlap in
every case (\cref{fig:cc-correlation}). The table does not say which estimator is closer to the truth, which only the simulator can: see
\cref{tab:cc-estacc}. It also does not say that the IRB correlation is wrong for the Board's series, because a yearly all-bank rate
averages the idiosyncratic noise of many portfolios and a fit on one calm window has no information about a regime that the window does
not contain.

\begin{figure}[htbp]
\centering
\includegraphics[width=0.95\textwidth]{credit-capital-correlation.pdf}
\caption{Asset correlation of the five loan types: the maximum-likelihood estimate with its 95 percent profile-likelihood interval, the
method-of-moments estimate with its 95 percent parametric-bootstrap interval of 500 draws, and the correlation of the IRB class of the
loan type (diamond), on the fit window 1985 to 2006 of the Federal Reserve charge-off rates; the horizontal axis is logarithmic.}
\label{fig:cc-correlation}
\end{figure}

\begin{figure}[htbp]
\centering
\includegraphics[width=0.95\textwidth]{credit-capital-profile.pdf}
\caption{Profile log-likelihood of the asset correlation minus its maximum for the five loan types, on the fit window 1985 to 2006
(1991 to 2006 for the two real estate series); the horizontal line is the 95 percent likelihood-ratio cutoff of 1.920729 below the
maximum, and the interval of each loan type is the stretch of its curve above the line.}
\label{fig:cc-profile}
\end{figure}

\begin{table}[htbp]
\centering
\small
\caption{Capital per unit of exposure at the 99.9\% level with an expected loss given default of 0.45: the IRB capital requirement at the regulatory correlation, the capital of the large homogeneous portfolio at the fitted correlation (both with the same maturity treatment), and their ratio, which isolates the effect of the correlation; the default probability is the fitted one. The last row repeats commercial real estate under the corporate curve, the class of general income-producing real estate, beside the high-volatility class of the row above.}
\label{tab:cc-irb}
\begin{tabular}{>{\raggedright\arraybackslash}p{2.8cm}>{\raggedright\arraybackslash}p{2.4cm}>{\raggedleft\arraybackslash}p{1.5cm}>{\raggedleft\arraybackslash}p{1.5cm}>{\raggedleft\arraybackslash}p{1.5cm}>{\raggedleft\arraybackslash}p{1.5cm}>{\raggedleft\arraybackslash}p{1.5cm}}
\toprule
Loan type & IRB class & $\rho_{\mathrm{IRB}}$ & $\hat\rho$ & $K_{\mathrm{IRB}}$ (\%) & $K_{\mathrm{LHP}}$ (\%) & Ratio \\
\midrule
Commercial and industrial & corporate & 0.196 & 0.064 & 7.14 & 2.40 & 0.34 \\
Commercial real estate & hvcre & 0.265 & 0.238 & 6.68 & 5.87 & 0.88 \\
Residential real estate & residential mortgage & 0.150 & 0.013 & 1.10 & 0.11 & 0.10 \\
Credit cards & qrre & 0.040 & 0.013 & 3.92 & 1.83 & 0.47 \\
Other consumer & other retail & 0.122 & 0.011 & 3.62 & 0.53 & 0.15 \\
Commercial real estate (corporate curve) & corporate & 0.216 & 0.238 & 5.23 & 5.87 & 1.12 \\
\bottomrule
\end{tabular}
\end{table}


\Cref{tab:cc-irb} puts the fitted correlation next to the IRB class. The IRB correlations are \ccRhoIrbRes{} for residential mortgages,
\ccRhoIrbCard{} for qualifying revolving exposures, \ccRhoIrbOth{} for other retail at the fitted default probability, \ccRhoIrbCi{} for
corporates and \ccRhoIrbCre{} for the commercial real estate class, and the capital of the large portfolio at the fitted correlation is
between \ccRatioMin{} and \ccRatioMax{} of the IRB capital: \ccRatioCre{} for commercial real estate, \ccRatioCard{} for credit cards,
\ccRatioCi{} for commercial and industrial loans, \ccRatioOth{} for other consumer and \ccRatioRes{} for residential real estate. Only the
commercial real estate fit is close to the regulatory value. The ratio is a ratio of two capital figures that differ only in the
correlation, so it measures how much of the IRB capital the loss history of the fit window supports, and says nothing yet about
2007 to 2012, which the next subsection tests. The regimes of \cref{tab:cc-regimes} change the capital by the Basel II scaling of
\ccRegScale{} and, for the Chilean regime, leave the risk weight of every class it accepts unchanged.

\subsubsection{The out-of-sample tail}

\begin{figure}[htbp]
\centering
\includegraphics[width=\textwidth]{credit-capital-tail.pdf}
\caption{Yearly net charge-off rates of the five loan types from 1985 to 2012 (percent per year; filled circles for the fit window,
open squares for the evaluation window 2007 to 2012; the vertical line marks the end of the fit window), with the median and the 95,
99 and 99.9 percent quantiles of the Vasicek distribution fitted on 1985 to 2006 as horizontal lines. Real estate series start in 1991.}
\label{fig:cc-tail}
\end{figure}

\begin{table}[htbp]
\centering
\small
\caption{Out-of-sample tail back-test: the number of evaluation years 2007 to 2012 whose loss rate exceeds the 95\%, 99\% and 99.9\% quantile of the distribution fitted on 1985 to 2006, the fitted 99\% quantile in percent, the upper binomial tail probability of the observed count at 99\% and 99.9\%, and the traffic light at the 99\% level (green above 0.05, amber above 0.01, red below). The expected counts are 0.3, 0.06 and 0.006.}
\label{tab:cc-tail}
\begin{tabular}{>{\raggedright\arraybackslash}p{2.6cm}>{\raggedleft\arraybackslash}p{1.4cm}>{\raggedleft\arraybackslash}p{1.3cm}>{\raggedleft\arraybackslash}p{1.3cm}>{\raggedleft\arraybackslash}p{1.3cm}>{\raggedleft\arraybackslash}p{1.5cm}>{\raggedleft\arraybackslash}p{1.5cm}>{\raggedright\arraybackslash}p{1.0cm}}
\toprule
Loan type & Fitted 99\% quantile (\%) & Years above 95\% & Years above 99\% & Years above 99.9\% & $p$ (99\%) & $p$ (99.9\%) & Light (99\%) \\
\midrule
Commercial and industrial & 3.34 & 1 of 6 & 0 of 6 & 0 of 6 & 1.000 & 1.000 & green \\
Commercial real estate & 4.42 & 2 of 6 & 0 of 6 & 0 of 6 & 1.000 & 1.000 & green \\
Residential real estate & 0.30 & 6 of 6 & 5 of 6 & 5 of 6 & \ensuremath{6.0\times10^{-10}} & \ensuremath{6.0\times10^{-15}} & red \\
Credit cards & 7.06 & 2 of 6 & 2 of 6 & 2 of 6 & 0.001 & \ensuremath{1.5\times10^{-5}} & red \\
Other consumer & 1.77 & 4 of 6 & 3 of 6 & 2 of 6 & \ensuremath{2.0\times10^{-5}} & \ensuremath{1.5\times10^{-5}} & red \\
\bottomrule
\end{tabular}
\end{table}


\Cref{fig:cc-tail} and \cref{tab:cc-tail} place the evaluation years in the distribution fitted before them. Under the fit, the expected
number of years above the 95, 99 and 99.9 percent quantiles out of \ccEvYears{} is \ccExpNF{}, \ccExpNN{} and \ccExpNNN{}. For residential real
estate \ccExcNFRes{} years exceed the 95 percent quantile of \ccQNFRes{} percent, \ccExcNNRes{} exceed the 99 percent quantile of
\ccQNNRes{} percent (the years \ccYrsNNRes{}) and \ccExcNNNRes{} exceed the 99.9 percent quantile of \ccQNNNRes{} percent, against a peak rate
of \ccPeakRateRes{} percent in \ccPeakYearRes{}; the upper binomial tail probability at 99 percent is \ccPNNRes{}. Credit cards
exceed their 99 percent quantile of \ccQNNCard{} percent in \ccExcNNCard{} years (\ccYrsNNCard{}, $p=\ccPNNCard$) and other consumer in
\ccExcNNOth{} years (\ccYrsNNOth{}, $p=\ccPNNOth$), while commercial and industrial and commercial real estate exceed it in none
(\ccExcNNCi{} and \ccExcNNCre{}); commercial real estate exceeds its 95 percent quantile in \ccExcNFCre{} years ($p=\ccPNFCre$). The
residential series is the extreme case: a model fitted on a window whose maximum rate is \ccFitMaxRateRes{} percent meets years of
\ccPeakRateRes{} percent, a regime change that no stationary Vasicek distribution produces, and the exceedance is as much a statement
about the stability of the default probability as about the correlation. The tail test has little power with six years, and what it
shows is therefore the direction of the failure on the three consumer and residential series, with probabilities far below the
traffic-light thresholds, together with the absence of failure on the two commercial series. It does not show which parameter of the model
is wrong.

\subsubsection{The models compared on the illustrative portfolios}

\begin{table}[htbp]
\centering
\small
\caption{Value at risk at the 99.9\% level of the illustrative portfolio of 100 million USD for each loan type, by model, in millions of USD: the expected loss (EL), the large homogeneous portfolio (LHP), the one-factor (G-1F) and two-sector (G-2F) Gaussian copulas, the Student-t copula with 5 degrees of freedom, CreditRisk+ (CR+), the large-portfolio quantile plus the granularity adjustment (LHP+GA), and importance sampling of the Gaussian copula (IS). All models share the fitted default probability and correlation.}
\label{tab:cc-var}
\begin{tabular}{>{\raggedright\arraybackslash}p{2.0cm}>{\raggedleft\arraybackslash}p{1.0cm}>{\raggedleft\arraybackslash}p{1.15cm}>{\raggedleft\arraybackslash}p{1.15cm}>{\raggedleft\arraybackslash}p{1.15cm}>{\raggedleft\arraybackslash}p{1.3cm}>{\raggedleft\arraybackslash}p{1.15cm}>{\raggedleft\arraybackslash}p{1.5cm}>{\raggedleft\arraybackslash}p{1.15cm}}
\toprule
Loan type & EL & LHP & G-1F & G-2F & Stud. $t$ & CR+ & LHP+GA & IS \\
\midrule
Commercial and industrial & 0.41 & 2.30 & 3.18 & 3.03 & 13.10 & 3.08 & 3.28 & 3.20 \\
Commercial real estate & 0.20 & 4.54 & 5.47 & 4.91 & 14.38 & 5.08 & 6.52 & 5.47 \\
Residential real estate & 0.06 & 0.18 & 0.41 & 0.41 & 5.98 & 0.41 & 0.33 & 0.41 \\
Credit cards & 1.90 & 3.74 & 3.87 & 3.61 & 15.60 & 3.83 & 3.83 & 3.87 \\
Other consumer & 0.44 & 0.97 & 1.12 & 1.04 & 11.20 & 1.10 & 1.05 & 1.12 \\
\bottomrule
\end{tabular}
\end{table}


\Cref{tab:cc-var} gives the 99.9 percent value at risk of each loan type's portfolio under each model, \cref{tab:cc-es} the expected shortfall,
and \cref{fig:cc-ladder} the value at risk relative to the large-portfolio value. The one-factor Gaussian copula is above the large
homogeneous portfolio by factors between \ccGoneOverLhpMin{} (credit cards, with \ccEffNamesCard{} effective names) and \ccGoneOverLhpMax{}
(residential real estate, where the default probability is so small that a few names default in the tail scenarios): the finite
portfolio matters more than the correlation for the low-default books. The two-sector Gaussian copula, with correlation \ccSectorCorr{} between
the sectors, is at \ccGmfOverGoneMin{} to \ccGmfOverGoneMax{} of the one-factor value, lower where the sectors diversify. CreditRisk+ with a
variance matched to the Vasicek default-rate variance lands at \ccCrpOverGoneMin{} to \ccCrpOverGoneMax{} of the Gaussian copula
(\cref{tab:cc-crp}), and the large-portfolio quantile plus the granularity adjustment at \ccGaOverGoneRes{} (residential real estate) to
\ccGaOverGoneCre{} (commercial real estate) of it. The Student-t copula with $\nu=\ccNu$ is \ccStuOverGoneMin{} to \ccStuOverGoneMax{} times the Gaussian
copula: its tail dependence, and not the correlation, sets the loss at 99.9 percent when the correlation is as small as the fits give, so
the choice of copula is the largest model risk of the comparison. The expected shortfall of the Student-t copula is \ccEsStuCard{} million USD for
credit cards against \ccEsGoneCard{} under the Gaussian copula. The comparison does not show that the Student-t copula is the better model:
$\nu=\ccNu$ is an input that is poorly identified from yearly data, and the exceedances above show only that the fitted Gaussian
distribution is too thin, not by how much.

\begin{figure}[htbp]
\centering
\includegraphics[width=0.95\textwidth]{credit-capital-ladder.pdf}
\caption{Value at risk at the 99.9 percent level of the illustrative portfolio of each loan type for each model, as a multiple of the
large-portfolio value (logarithmic scale); error bars are plus and minus 1.96 batch-means standard errors for the sampled models, and the
analytic models carry none. The Student-t copula with 5 degrees of freedom lies above every other model for all five loan types.}
\label{fig:cc-ladder}
\end{figure}

\begin{table}[htbp]
\centering
\small
\caption{Expected shortfall at the 99.9\% level of the same portfolios and models as the value-at-risk table, with the same column labels, in millions of USD; the granularity adjustment is a quantile add-on and defines no expected shortfall, so its column is absent. The Student-t copula carries by far the heaviest tail beyond the quantile.}
\label{tab:cc-es}
\begin{tabular}{>{\raggedright\arraybackslash}p{2.8cm}>{\raggedleft\arraybackslash}p{1.6cm}>{\raggedleft\arraybackslash}p{1.6cm}>{\raggedleft\arraybackslash}p{1.6cm}>{\raggedleft\arraybackslash}p{1.6cm}>{\raggedleft\arraybackslash}p{1.6cm}>{\raggedleft\arraybackslash}p{1.6cm}}
\toprule
Loan type & LHP & G-1F & G-2F & Stud. $t$ & CR+ & IS \\
\midrule
Commercial and industrial & 2.68 & 3.68 & 3.50 & 16.00 & 3.51 & 3.71 \\
Commercial real estate & 6.01 & 7.08 & 6.21 & 18.97 & 6.16 & 7.06 \\
Residential real estate & 0.20 & 0.46 & 0.45 & 8.23 & 0.46 & 0.46 \\
Credit cards & 3.96 & 4.11 & 3.81 & 17.13 & 4.05 & 4.11 \\
Other consumer & 1.04 & 1.21 & 1.12 & 13.23 & 1.18 & 1.21 \\
\bottomrule
\end{tabular}
\end{table}


\begin{figure}[htbp]
\centering
\includegraphics[width=\textwidth]{credit-capital-exceedance.pdf}
\caption{Probability that the loss of the illustrative portfolio exceeds a given amount (million USD, logarithmic probability axis) for
the large-portfolio formula, the one-factor and two-sector Gaussian copulas, the Student-t copula and CreditRisk+, one panel per loan type;
the horizontal line marks the probability 0.001 of the 99.9 percent level. Sampled curves are cut where the exceedance probability falls
below $10^{-5}$.}
\label{fig:cc-exceedance}
\end{figure}

\Cref{fig:cc-exceedance} shows why: the Gaussian curves and the large-portfolio curve are close, while the Student-t curve crosses the 0.001
line at a loss several times larger. Importance sampling reproduces the plain Monte Carlo value at risk (ratio \ccIsOverGoneMin{} to
\ccIsOverGoneMax{} of the plain value across the five loan types, within one smallest-name loss on the lattice of the portfolio) with a variance
ratio of \ccIsVarRatioMin{} to \ccIsVarRatioMax{} for the value at risk and \ccIsEsRatioMin{} to \ccIsEsRatioMax{} for the expected shortfall
(\cref{tab:cc-mcerr}). The effective sample share of the weights is small, from \ccIsEssCard{} percent for credit cards to \ccIsEssRes{} percent for
residential real estate, so a few scenarios carry most of the weight: the variance ratio is a statement about the batch-means error and must be read
with that share, and the number of scenarios is raised rather than trusted to the mean of the weights.

\begin{table}[htbp]
\centering
\small
\caption{Monte Carlo error of the 99.9\% value at risk, as the batch-means standard error in thousands of USD, for the three copula samples of $10^6$ scenarios and the importance-sampling sample of $10^5$ scenarios; variance ratios (VR) of importance sampling against plain Monte Carlo at equal sample size for the value at risk and the expected shortfall; and the effective sample share (ESS) of the importance weights in percent; labels as in the value-at-risk table.}
\label{tab:cc-mcerr}
\begin{tabular}{>{\raggedright\arraybackslash}p{2.4cm}>{\raggedleft\arraybackslash}p{1.4cm}>{\raggedleft\arraybackslash}p{1.4cm}>{\raggedleft\arraybackslash}p{1.4cm}>{\raggedleft\arraybackslash}p{1.4cm}>{\raggedleft\arraybackslash}p{1.4cm}>{\raggedleft\arraybackslash}p{1.4cm}>{\raggedleft\arraybackslash}p{1.4cm}}
\toprule
Loan type & SE G-1F & SE G-2F & SE Stud. $t$ & SE IS & VR VaR & VR ES & ESS (\%) \\
\midrule
Commercial and industrial & 15.2 & 15.0 & 111.6 & 2.4 & 389 & 644 & 0.017 \\
Commercial real estate & 52.6 & 41.9 & 132.4 & 8.7 & 369 & 1217 & 0.077 \\
Residential real estate & 1.4 & 1.1 & 88.4 & 0.5 & 93 & 995 & 1.708 \\
Credit cards & 6.4 & 7.0 & 57.9 & 1.8 & 120 & 237 & 0.005 \\
Other consumer & 1.8 & 2.7 & 68.4 & 0.7 & 60 & 506 & 0.006 \\
\bottomrule
\end{tabular}
\end{table}


\begin{table}[htbp]
\centering
\small
\caption{CreditRisk+ with one gamma sector whose variance $\sigma^2$ matches the Vasicek default-rate variance of the fit: the probability-mass grid length in loss units of 1,000 USD, the probability mass left beyond the grid (a negative value is a rounding excess of the summed mass over one), the expected loss recovered from the mass function against the portfolio's own, and the 99.9\% value at risk relative to the one-factor Gaussian copula.}
\label{tab:cc-crp}
\begin{tabular}{>{\raggedright\arraybackslash}p{2.5cm}>{\raggedleft\arraybackslash}p{1.6cm}>{\raggedleft\arraybackslash}p{1.6cm}>{\raggedleft\arraybackslash}p{1.6cm}>{\raggedleft\arraybackslash}p{1.6cm}>{\raggedleft\arraybackslash}p{1.6cm}>{\raggedleft\arraybackslash}p{1.6cm}}
\toprule
Loan type & $\sigma^2$ & Grid $n_{\max}$ & Mass beyond grid & EL of the pmf (M) & EL of the portfolio (M) & VaR / Gaussian 1F \\
\midrule
Commercial and industrial & 0.5521 & 12,728 & \ensuremath{-1.5\times10^{-13}} & 0.410 & 0.410 & 0.969 \\
Commercial real estate & 4.5370 & 15,983 & \ensuremath{6.3\times10^{-8}} & 0.196 & 0.196 & 0.928 \\
Residential real estate & 0.1464 & 4,000 & \ensuremath{2.0\times10^{-13}} & 0.063 & 0.063 & 1.000 \\
Credit cards & 0.0590 & 8,506 & \ensuremath{4.3\times10^{-12}} & 1.903 & 1.903 & 0.988 \\
Other consumer & 0.0805 & 4,000 & \ensuremath{-6.0\times10^{-13}} & 0.439 & 0.439 & 0.976 \\
\bottomrule
\end{tabular}
\end{table}


\subsubsection{Granularity and allocation}

\begin{table}[htbp]
\centering
\small
\caption{Granularity adjustment of Gordy and L\"utkebohmert for the illustrative portfolios: the portfolio capital $K^{*}$ from the IRB formula at the regulatory correlation, the adjustment as a share of exposure at the 99\% and 99.9\% levels, the adjustment relative to the large-portfolio quantile, the gap of the one-factor Gaussian copula above the large-portfolio quantile on the same scale, and the weight on the systematic factor implied by the IRB capital of one name.}
\label{tab:cc-ga}
\begin{tabular}{>{\raggedright\arraybackslash}p{2.5cm}>{\raggedleft\arraybackslash}p{1.7cm}>{\raggedleft\arraybackslash}p{1.7cm}>{\raggedleft\arraybackslash}p{1.7cm}>{\raggedleft\arraybackslash}p{1.7cm}>{\raggedleft\arraybackslash}p{1.7cm}>{\raggedleft\arraybackslash}p{1.7cm}}
\toprule
Loan type & $K^{*}$ (\%) & GA 99\% (\% EAD) & GA 99.9\% (\% EAD) & GA / LHP VaR (\%) & Copula gap / LHP VaR (\%) & Implied $\omega$ \\
\midrule
Commercial and industrial & 6.90 & 0.482 & 0.983 & 42.8 & 38.2 & 1.055 \\
Commercial real estate & 6.45 & 0.966 & 1.982 & 43.7 & 20.5 & 2.062 \\
Residential real estate & 1.09 & 0.077 & 0.156 & 87.8 & 127.8 & 1.064 \\
Credit cards & 3.87 & 0.050 & 0.095 & 2.5 & 3.6 & 0.125 \\
Other consumer & 3.51 & 0.043 & 0.087 & 9.0 & 16.3 & 0.500 \\
\bottomrule
\end{tabular}
\end{table}


\Cref{tab:cc-ga} gives the adjustment: \ccGaTopRes{} percent of the exposure for residential real estate, \ccGaTopCard{} for credit cards,
\ccGaTopOth{} for other consumer, \ccGaTopCi{} for commercial and industrial and \ccGaTopCre{} for commercial real estate at the 99.9 percent
level, which is \ccGaShareMin{} to \ccGaShareMax{} percent of the large-portfolio value at risk across the loan types (\ccGaShareCard{}
percent for credit cards, \ccGaShareRes{} percent for residential real estate). The multiplier is $\delta=\ccDelta$ at the 99.9 percent level and $\ccDeltaLow$ at 99 percent ($\alpha_q=\ccAlphaQ$ and
$\ccAlphaQLow$). The adjustment tracks the direction of the finite-portfolio gap of the Gaussian copula (the large-portfolio value plus
the adjustment is at \ccGaOverGoneRes{} to \ccGaOverGoneCre{} of the copula value) but it is a first-order formula and does not reproduce the
gap in magnitude. For the three loan types in which the IRB capital of one name implies a weight on the systematic factor above one
($\omega=\ccOmegaCi$, $\ccOmegaCre$ and $\ccOmegaRes$ for commercial and industrial, commercial real estate and residential real estate, last
column of \cref{tab:cc-ga}) the CreditRisk+ reference is not defined and the document says so; for credit cards ($\omega=\ccOmegaCard$) and
other consumer ($\omega=\ccOmegaOth$) it is, and \cref{tab:cc-ga-n} and \cref{fig:cc-granularity} set the adjustment against it for portfolios of
10 to 2,000 identical names. The adjustment is above the reference at every size, by factors of \ccGaRatioMin{} to \ccGaRatioMax{},
and falls as $1/n$, as it must. The reference includes no variance of the loss given default, which accounts for part of the excess, and
it moves in whole-name steps (the reference of the other-consumer mix is the same for 10 and 20 names). The table does not show that the
adjustment is accurate for a real portfolio, whose names are not identical.

\begin{table}[htbp]
\centering
\small
\caption{Granularity adjustment against its reference at the 99.9\% level for portfolios of $n$ identical names of the credit-card and other-consumer mixes, in percent of exposure: the adjustment of the formula and the gap of the CreditRisk+ recursion of the same single-factor model above the large-portfolio quantile (the reference), with their ratio. The ratio exceeds one at every size, and the adjustment scales as $1/n$.}
\label{tab:cc-ga-n}
\begin{tabular}{>{\raggedleft\arraybackslash}p{1.1cm}>{\raggedleft\arraybackslash}p{1.2cm}>{\raggedleft\arraybackslash}p{1.55cm}>{\raggedleft\arraybackslash}p{1.55cm}>{\raggedleft\arraybackslash}p{1.55cm}>{\raggedleft\arraybackslash}p{1.55cm}>{\raggedleft\arraybackslash}p{1.55cm}>{\raggedleft\arraybackslash}p{1.55cm}}
\toprule
Names & HHI & GA cards & Recursion cards & Ratio & GA other & Recursion other & Ratio \\
\midrule
10 & 0.1000 & 18.158 & 12.180 & 1.49 & 12.981 & 4.939 & 2.63 \\
20 & 0.0500 & 9.079 & 5.430 & 1.67 & 6.491 & 4.939 & 1.31 \\
50 & 0.0200 & 3.632 & 2.280 & 1.59 & 2.596 & 1.339 & 1.94 \\
100 & 0.0100 & 1.816 & 1.380 & 1.32 & 1.298 & 0.889 & 1.46 \\
200 & 0.0050 & 0.908 & 0.705 & 1.29 & 0.649 & 0.439 & 1.48 \\
500 & 0.0020 & 0.363 & 0.300 & 1.21 & 0.260 & 0.169 & 1.54 \\
1,000 & 0.0010 & 0.182 & 0.120 & 1.51 & 0.130 & 0.079 & 1.64 \\
2,000 & 0.0005 & 0.091 & 0.075 & 1.20 & 0.065 & 0.056 & 1.15 \\
\bottomrule
\end{tabular}
\end{table}


\begin{figure}[htbp]
\centering
\includegraphics[width=0.95\textwidth]{credit-capital-granularity.pdf}
\caption{Gap of the 99.9 percent value at risk of a portfolio of $n$ identical names above the large-portfolio value, in percent of
exposure, against $n$ (logarithmic axes): the granularity adjustment of the formula and the CreditRisk+ recursion of the same single-factor
model, for the credit-card mix (left) and the other-consumer mix (right).}
\label{fig:cc-granularity}
\end{figure}

\begin{table}[htbp]
\centering
\small
\caption{Euler capital allocation of the one-factor Gaussian copula at the 99.9\% level by group of names: the share of exposure, the share of the expected-shortfall contribution, the share of the kernel-smoothed value-at-risk contribution, and the share of the asymptotic single-factor contribution of Emmer and Tasche; shares of each column sum to 100\%.}
\label{tab:cc-euler}
\begin{tabular}{>{\raggedright\arraybackslash}p{3.2cm}>{\raggedleft\arraybackslash}p{1.2cm}>{\raggedleft\arraybackslash}p{2.0cm}>{\raggedleft\arraybackslash}p{2.0cm}>{\raggedleft\arraybackslash}p{2.0cm}>{\raggedleft\arraybackslash}p{2.0cm}}
\toprule
Loan type & Group & Exposure (\%) & ES share (\%) & VaR share, kernel (\%) & Asymptotic share (\%) \\
\midrule
Commercial and industrial & large & 40.0 & 55.5 & 55.7 & 25.5 \\
 & mid & 40.0 & 29.2 & 29.2 & 44.2 \\
 & small & 20.0 & 15.3 & 15.1 & 30.3 \\
Commercial real estate & large & 40.0 & 44.1 & 46.3 & 27.2 \\
 & mid & 40.0 & 36.5 & 34.6 & 44.0 \\
 & small & 20.0 & 19.3 & 19.2 & 28.8 \\
Residential real estate & large & 20.0 & 43.7 & 46.4 & 9.8 \\
 & mid & 40.0 & 41.5 & 41.3 & 36.9 \\
 & small & 40.0 & 14.8 & 12.2 & 53.3 \\
Credit cards & large & 10.0 & 4.9 & 5.0 & 4.7 \\
 & mid & 40.0 & 37.1 & 37.2 & 34.3 \\
 & small & 50.0 & 58.0 & 57.6 & 61.0 \\
Other consumer & large & 20.0 & 17.7 & 17.1 & 9.9 \\
 & mid & 40.0 & 39.4 & 39.3 & 36.9 \\
 & small & 40.0 & 42.9 & 43.5 & 53.2 \\
\bottomrule
\end{tabular}
\end{table}


\Cref{tab:cc-euler} allocates the 99.9 percent capital of the one-factor Gaussian copula to the three groups. The expected-shortfall
contributions sum to the total to a relative error of at most \ccEulerEsWorst{} over the \ccEulerRows{} additivity checks of the five
portfolios (three copula samples, two levels, two measures), and the kernel value-at-risk contributions to at most \ccEulerVarWorst{}, inside
the tolerance of \ccEulerTolVar{}. The two measures agree with each other, within \ccEulerShareGap{} percentage points, in every group. The asymptotic
single-factor contributions do not: for commercial and industrial loans the large names carry \ccEulerLargeEsCi{} percent of the expected-shortfall
contribution against \ccEulerLargeAsyCi{} percent of the asymptotic one, because the asymptotic contribution is portfolio-invariant and
proportional to each name's own exposure and conditional default probability, so it cannot see that \ccLargeNCi{} names of \ccLargeEadCi{} million
USD make a lumpy book. For the consumer
books the shares are close because the groups are many small names. The table is for the one-factor Gaussian copula only; the multi-factor
and Student-t rungs have no closed-form anchor.

\subsubsection{The validation tests}

\begin{table}[htbp]
\centering
\small
\caption{Validation tests of the portfolio models on each loan type: the convergence of the Monte Carlo value at risk of 5,000 identical names to the large-portfolio quantile ($z$ statistic, two-sided $p$-value, light); the Kolmogorov distance $D$ between the CreditRisk+ recursion and a sample of the same model, with its $p$-value and light; and the worst relative gap between the sum of the Euler contributions and the total over the three copula samples and both levels.}
\label{tab:cc-battery}
\begin{tabular}{>{\raggedright\arraybackslash}p{2.3cm}>{\raggedleft\arraybackslash}p{1.2cm}>{\raggedleft\arraybackslash}p{1.2cm}>{\raggedleft\arraybackslash}p{1.2cm}>{\raggedleft\arraybackslash}p{1.2cm}>{\raggedleft\arraybackslash}p{1.2cm}>{\raggedleft\arraybackslash}p{1.2cm}>{\raggedleft\arraybackslash}p{1.2cm}>{\raggedleft\arraybackslash}p{1.2cm}}
\toprule
Loan type & MC $z$ & MC $p$ & Light & CR+ $D$ & CR+ $p$ & Light & VaR gap & ES gap \\
\midrule
Commercial and industrial & $-0.54$ & 0.587 & green & 0.00126 & 0.554 & green & 0.0015 & \ensuremath{2.1\times10^{-14}} \\
Commercial real estate & $-1.51$ & 0.131 & green & 0.00069 & 0.992 & green & 0.0037 & \ensuremath{1.7\times10^{-14}} \\
Residential real estate & $23.38$ & \ensuremath{6.3\times10^{-121}} & red & 0.00086 & 0.930 & green & 0.0016 & \ensuremath{8.1\times10^{-15}} \\
Credit cards & $1.27$ & 0.205 & green & 0.00084 & 0.940 & green & 0.0018 & \ensuremath{1.3\times10^{-15}} \\
Other consumer & $7.73$ & \ensuremath{1.1\times10^{-14}} & red & 0.00100 & 0.819 & green & 0.0013 & \ensuremath{5.2\times10^{-15}} \\
\bottomrule
\end{tabular}
\end{table}


\Cref{tab:cc-battery} reports the tests on each portfolio. The Monte Carlo value at risk of 5,000 identical names agrees with the
large-portfolio quantile for commercial and industrial ($z=\ccMcZCi$, $p=\ccMcPCi$), commercial real estate ($z=\ccMcZCre$, $p=\ccMcPCre$)
and credit cards ($z=\ccMcZCard$, $p=\ccMcPCard$), and does not for residential real estate ($z=\ccMcZRes$) and other consumer
($z=\ccMcZOth$): their default probabilities are low enough that 5,000 names still carry only a handful of defaults in the tail scenarios, so
the sample sits on a lattice above the continuous quantile, and the check reports the gap rather than hide it. The CreditRisk+ recursion matches a
Monte Carlo sample of its own model in every case, with $D$ between \ccCrpDCre{} and \ccCrpDCi{} and $p$ between \ccCrpPCi{} and
\ccCrpPCre{}. The tail back-test lights are \ccBtLightCi{} for commercial and industrial, \ccBtLightCre{} for commercial real estate and
\ccBtLightRes{} for residential real estate, credit cards and other consumer. A green light does not say that the model is adequate, and none of
these tests validates the choice of the model, the portfolio or the dependence assumption.

\subsubsection{Size and power of the correlation estimators}

\begin{table}[htbp]
\centering
\small
\caption{Rejection rate of $H_0:\rho=0.12$ at the 5\% level on the simulator with known correlation and a known default probability of 2\%, for the maximum-likelihood and method-of-moments tests with series of 22 and 40 years: 400 replications per cell, the Wilson 95\% interval in square brackets. The row at 0.12 is the size of the test; every other row is its power.}
\label{tab:cc-sizepower}
\begin{tabular}{>{\raggedleft\arraybackslash}p{0.9cm}>{\raggedright\arraybackslash}p{1.0cm}>{\raggedleft\arraybackslash}p{2.75cm}>{\raggedleft\arraybackslash}p{2.75cm}>{\raggedleft\arraybackslash}p{2.75cm}>{\raggedleft\arraybackslash}p{2.75cm}}
\toprule
$\rho$ & Role & ML, 22 y & Moments, 22 y & ML, 40 y & Moments, 40 y \\
\midrule
0.02 & power & 1.000 [0.990, 1.000] & 1.000 [0.990, 1.000] & 1.000 [0.990, 1.000] & 1.000 [0.990, 1.000] \\
0.04 & power & 0.990 [0.975, 0.996] & 0.985 [0.968, 0.993] & 1.000 [0.990, 1.000] & 0.993 [0.978, 0.997] \\
0.06 & power & 0.740 [0.695, 0.781] & 0.812 [0.771, 0.848] & 0.938 [0.909, 0.957] & 0.897 [0.864, 0.924] \\
0.08 & power & 0.318 [0.274, 0.365] & 0.585 [0.536, 0.632] & 0.550 [0.501, 0.598] & 0.615 [0.566, 0.661] \\
0.10 & power & 0.152 [0.121, 0.191] & 0.350 [0.305, 0.398] & 0.147 [0.116, 0.186] & 0.297 [0.255, 0.344] \\
0.12 & size & 0.052 [0.035, 0.079] & 0.185 [0.150, 0.226] & 0.060 [0.041, 0.088] & 0.130 [0.101, 0.167] \\
0.14 & power & 0.075 [0.053, 0.105] & 0.122 [0.094, 0.158] & 0.107 [0.081, 0.142] & 0.090 [0.066, 0.122] \\
0.16 & power & 0.207 [0.171, 0.250] & 0.065 [0.045, 0.094] & 0.270 [0.229, 0.316] & 0.085 [0.061, 0.116] \\
0.18 & power & 0.287 [0.245, 0.334] & 0.048 [0.031, 0.073] & 0.560 [0.511, 0.608] & 0.107 [0.081, 0.142] \\
0.20 & power & 0.495 [0.446, 0.544] & 0.087 [0.064, 0.119] & 0.775 [0.732, 0.813] & 0.145 [0.114, 0.183] \\
0.22 & power & 0.632 [0.584, 0.678] & 0.068 [0.047, 0.096] & 0.875 [0.839, 0.904] & 0.220 [0.182, 0.263] \\
0.24 & power & 0.725 [0.679, 0.766] & 0.085 [0.061, 0.116] & 0.968 [0.945, 0.981] & 0.305 [0.262, 0.352] \\
\bottomrule
\end{tabular}
\end{table}


\begin{figure}[htbp]
\centering
\includegraphics[width=\textwidth]{credit-capital-sizepower.pdf}
\caption{Rejection rate of the test of an asset correlation of 0.12 at the 5 percent level against the known correlation of the simulator
(default probability 2 percent), for the maximum-likelihood test with its profile interval and the method-of-moments test with its bootstrap
interval; series of 22 years (left) and 40 years (right), 400 replications per point, bars are Wilson 95 percent intervals. The horizontal
line is the nominal 5 percent and the dotted vertical line the null correlation.}
\label{fig:cc-sizepower}
\end{figure}

The simulator draws one standard normal factor per year, converts it to a default fraction with \cref{eq:cc-condpd} at a known default probability of
\ccSimPd{} percent, and fits the series with both estimators on a grid of \ccNGrid{} correlations from \ccRhoGridLo{} to \ccRhoGridHi{}, with series of
\ccSerShort{} and \ccSerLong{} years and \ccSimRep{} replications per cell, each with its own seeded random stream (seed \ccSimSeed{}); the null is
$\rho=\ccNullRho$ at the \ccAlpha{} percent level, rejected when the null lies outside the estimator's 95 percent interval, and
no replication failed to identify a correlation (\ccNFailedShort{} failures with 22 years and \ccNFailedLong{} with 40). \Cref{tab:cc-sizepower} and \cref{fig:cc-sizepower} give the rejection rate. At the null
correlation the maximum-likelihood test rejects \ccSizeMlShort{} of the time with 22 years (Wilson interval \ccSizeMlShortLo{} to \ccSizeMlShortHi{})
and \ccSizeMlLong{} with 40 years, close to the nominal 0.05; the method-of-moments test rejects \ccSizeMoShort{} (\ccSizeMoShortLo{} to
\ccSizeMoShortHi{}) and \ccSizeMoLong{}, which is \ccSizeMoRatioShort{} and \ccSizeMoRatioLong{} times the nominal rate. The power of the maximum-likelihood test is \ccPowMlShortFour{} at a
correlation of 0.04 and \ccPowMlShortEight{} at 0.08 with 22 years, against \ccPowMlLongFour{} and \ccPowMlLongEight{} with 40 years; at 0.24 it is
\ccPowMlShortTwoFour{} and \ccPowMlLongTwoFour{}, whereas the moments test has power \ccPowMoShortTwoFour{} and \ccPowMoLongTwoFour{} there: its estimates are biased
downward (the most negative bias is $\ccBiasMoWorstShort$ with 22 years, against $\ccBiasMlWorstShort$ for maximum likelihood) and more dispersed
(\cref{tab:cc-estacc}), so its interval keeps containing the null correlation when the true one lies above it.

\begin{table}[htbp]
\centering
\small
\caption{Accuracy of the two correlation estimators on the simulator with known correlation: root-mean-square error of the estimate and empirical coverage of its nominal 95\% interval, for series of 22 and 40 years, over 400 replications per row; the coverage has a Monte Carlo standard error of about 0.011 near 0.95.}
\label{tab:cc-estacc}
\begin{tabular}{>{\raggedleft\arraybackslash}p{0.8cm}>{\raggedleft\arraybackslash}p{1.3cm}>{\raggedleft\arraybackslash}p{1.3cm}>{\raggedleft\arraybackslash}p{1.3cm}>{\raggedleft\arraybackslash}p{1.3cm}>{\raggedleft\arraybackslash}p{1.3cm}>{\raggedleft\arraybackslash}p{1.3cm}>{\raggedleft\arraybackslash}p{1.3cm}>{\raggedleft\arraybackslash}p{1.3cm}}
\toprule
$\rho$ & ML RMSE 22 & ML cov. 22 & Mom. RMSE 22 & Mom. cov. 22 & ML RMSE 40 & ML cov. 40 & Mom. RMSE 40 & Mom. cov. 40 \\
\midrule
0.02 & 0.006 & 0.912 & 0.007 & 0.823 & 0.005 & 0.943 & 0.005 & 0.858 \\
0.04 & 0.011 & 0.932 & 0.013 & 0.843 & 0.009 & 0.940 & 0.011 & 0.887 \\
0.06 & 0.016 & 0.953 & 0.022 & 0.823 & 0.012 & 0.935 & 0.016 & 0.905 \\
0.08 & 0.022 & 0.938 & 0.027 & 0.820 & 0.017 & 0.932 & 0.023 & 0.838 \\
0.10 & 0.027 & 0.930 & 0.033 & 0.823 & 0.019 & 0.953 & 0.029 & 0.858 \\
0.12 & 0.030 & 0.948 & 0.039 & 0.815 & 0.024 & 0.940 & 0.033 & 0.870 \\
0.14 & 0.036 & 0.945 & 0.048 & 0.777 & 0.027 & 0.938 & 0.041 & 0.850 \\
0.16 & 0.040 & 0.953 & 0.057 & 0.790 & 0.030 & 0.950 & 0.051 & 0.815 \\
0.18 & 0.045 & 0.930 & 0.064 & 0.738 & 0.032 & 0.953 & 0.056 & 0.815 \\
0.20 & 0.049 & 0.943 & 0.079 & 0.745 & 0.035 & 0.953 & 0.057 & 0.863 \\
0.22 & 0.052 & 0.955 & 0.075 & 0.762 & 0.039 & 0.943 & 0.068 & 0.807 \\
0.24 & 0.057 & 0.940 & 0.094 & 0.655 & 0.038 & 0.950 & 0.070 & 0.820 \\
\bottomrule
\end{tabular}
\end{table}


\Cref{tab:cc-estacc} shows the accuracy behind those rates. The maximum-likelihood interval has an empirical coverage between \ccCovMlMinShort{} and
\ccCovMlMaxShort{} with 22 years and between \ccCovMlMinLong{} and \ccCovMlMaxLong{} with 40 years, around the nominal 0.95 with a Monte Carlo standard
error of about 0.011, while the bootstrap interval of the moments estimator covers only \ccCovMoMinShort{} to \ccCovMoMaxShort{} and \ccCovMoMinLong{} to
\ccCovMoMaxLong{}. Two consequences follow for the fits above. The maximum-likelihood intervals of \cref{tab:cc-fit} can be read at close to their nominal
level, when the simulator's assumptions (independent years, a known default probability, no regime change) hold; the moments intervals are too narrow, so
the agreement of the two estimators in \cref{fig:cc-correlation} is weaker evidence than the overlap suggests. The simulator does not show what happens when
the default probability itself moves, which is the situation the residential series presents, and the results hold for its generator and series lengths only.
\emergencystretch=0pt
